\documentclass[12pt,a4paper]{report}

\usepackage[T1]{fontenc}
\usepackage[utf8]{inputenc}
\usepackage{lmodern}
\usepackage[margin=2.35cm]{geometry}
\usepackage{setspace}
\usepackage{amsmath,amssymb,amsthm,mathtools}
\usepackage{booktabs,tabularx,longtable,multirow,array}
\usepackage{graphicx}
\usepackage{xcolor}
\usepackage{tikz}
\usetikzlibrary{arrows.meta,positioning,shapes.geometric,fit,calc,decorations.pathreplacing}
\usepackage{pgfplots}
\pgfplotsset{compat=1.17}
\usepackage{algorithm}
\usepackage{algpseudocode}
\usepackage{enumitem}
\usepackage{caption}
\usepackage{hyperref}
\hypersetup{colorlinks=true,linkcolor=black,citecolor=black,urlcolor=blue}

\onehalfspacing
\setstretch{1.45}
\setlength{\parskip}{0.35em}
\setlength{\parindent}{1.25em}
\raggedbottom

\newtheorem{definition}{Definition}[chapter]
\newtheorem{proposition}{Proposition}[chapter]
\newtheorem{remark}{Remark}[chapter]

\newcommand{\cmax}{C_{\max}}
\newcommand{\R}{\mathbb{R}}
\newcommand{\E}{\mathbb{E}}
\newcommand{\Prob}{\mathbb{P}}
\newcommand{\1}{\mathbf{1}}

\title{Research for Solving Multi-Target Scheduling of School Bag Factory Production by Using Resource-Coupled Graph Reinforcement Learning and Status-Synchronized Digital Twin Coordination}
\author{Computer Science Thesis\\System Design and Optimization}
\date{October 2026}

\begin{document}
\maketitle

\begin{abstract}
\noindent
A school bag factory is a flexible job shop with a particular kind of coupling. Cutting, embroidery, sewing, hardware assembly, inspection, and packing can each run on more than one machine, but an operation is not free to start merely because a machine is idle. It also needs a qualified operator, and often a shared token such as a cutting die or a colour-matched thread lot. While the schedule is running, the machines are not constant. Cutter wear, needle temperature, motor current, and thread-break counts change the time and the energy of the next cycle, and they change the chance that the cycle will be interrupted by a breakdown. The planner therefore has several targets at once: a short makespan, small tardiness against order due dates, low energy, balanced workload, and few failures. A rule that is excellent for one target is routinely poor for another.

This thesis studies that planning problem and proposes two technologies that are designed to work as one system. The first, called resource-coupled heterogeneous graph attention scheduling (RCH-GAT), represents the shop as a graph of operations, machines, operators, and shared resource tokens. A preference-conditioned attention policy chooses, at each decision epoch, an operation together with a machine and a resource bundle. One policy serves a family of target weights, so the same system can be asked to protect due dates in one shift and energy in the next. The second technology, called the status-synchronized digital twin with chance-constrained rescheduling (SSDT-CC), keeps a live degradation state for every machine. Processing times become functions of that state. A chance constraint refuses an assignment whose predicted failure probability exceeds a set tolerance, and inserts a short preventive maintenance instead. Rescheduling is driven by the twin, not by a fixed clock: the policy is called again when health, a breakdown, or a newly released rush order makes the current plan unreliable.

The first chapter reviews the two bodies of work these technologies come from, and sets the factory problem in a form that those bodies of work do not quite cover. The second chapter gives the models, the formulas, the decision procedures, and worked numerical steps. The third chapter describes the system that was built around the factory and reports a controlled estimation on a school-bag shop simulator. The estimation program is part of the thesis. On the nominal order book the full system improves the joint cost relative to first-in-first-out, shortest processing time, earliest due date, and apparent tardiness cost, and the two ablations (resource features removed, status features removed) each give back part of that gain. The results are estimates of this implementation, not figures copied from another factory.
\end{abstract}

\tableofcontents
\listoffigures
\listoftables

\chapter*{Notation}
\addcontentsline{toc}{chapter}{Notation}
\noindent
The symbols below are used with a stable meaning. Local symbols are defined where they appear.

\begin{longtable}{@{}p{0.22\textwidth}p{0.70\textwidth}@{}}
\toprule
Symbol & Meaning \\
\midrule
\endhead
$J$, $M$, $K$, $Q$ & Sets of jobs (orders), machines, operators, and shared resource tokens. \\
$O_{ij}$ & Operation $i$ of job $j$. \\
$G_{ij}$ & Set of machine groups that can perform $O_{ij}$. \\
$p_{ijko}(h)$ & Processing time of $O_{ij}$ on machine $k$ with operator $o$ at health $h$. \\
$s_{ij k}$ & Sequence-dependent setup on machine $k$ before $O_{ij}$. \\
$C_j$, $d_j$, $r_j$ & Completion time, due date, and release time of job $j$. \\
$\cmax$ & Makespan, $\max_j C_j$. \\
$T$ & Total tardiness, $\sum_j \max(0,C_j-d_j)$. \\
$E$ & Total energy of processing, setup, maintenance, and expected repair. \\
$B$ & Workload imbalance across machines. \\
$\Phi$ & Expected number of breakdowns. \\
$\omega$ & Preference vector over the targets. Its components are non-negative and sum to one. \\
$h_k(t)$ & Degradation state of machine $k$ at time $t$. \\
$\varepsilon$ & Chance-constraint tolerance on the failure probability of an assignment. \\
$\mathcal{G}$ & Heterogeneous shop graph. \\
$\alpha_{uv}^{(r)}$ & Attention weight on an edge of relation $r$. \\
$\pi_\theta(a\mid s,\omega)$ & Scheduling policy. \\
\bottomrule
\end{longtable}

\chapter{Research of Previous Work}
\label{ch:previous}

A school bag looks like a simple product. The shop that makes it is not simple. One order asks for a navy backpack in three sizes, the next for a red trolley bag with a short due date, and a third for a khaki sports sack that must share the only khaki thread lot with work already on a machine. The auto cutter is fast and already worn. The second cutter is slow and healthy. Three sewing operators do not all have the same skill on the bartack machine. Embroidery has a single head. Inspection and packing share people with the hardware bench. A planner who treats this shop as a list of jobs and a list of machines will build a schedule that looks feasible on paper and then loses the afternoon to a colour change, a missing operator, or a cutter that was asked to run past a sensible wear limit.

The scheduling literature already contains powerful answers to pieces of this picture. Flexible job shops, dispatching rules, evolutionary multi-objective search, and, more recently, graph neural networks trained by reinforcement learning, all decide which operation should run on which machine. A second literature, built around cyber-physical systems, digital twins, and remaining-useful-life models, decides what a machine's present health means for the next hour. Textile research has started to treat operators as a second resource. Almost none of these lines decides the school-bag problem as it actually arrives: several targets, a resource bundle around every operation, and a machine state that is allowed to change the plan.

This chapter reviews those lines in the order the thesis uses them. Section~\ref{sec:tech1-review} reviews scheduling technology, from classical job shops to heterogeneous graph policies, and states what that technology still does not represent. Section~\ref{sec:tech2-review} reviews machine-status technology, from degradation models to digital twins that influence a schedule, and states where the loop is still open. Section~\ref{sec:problem-set} puts the two gaps on the same factory, writes the problem the rest of the thesis solves, and names the two technologies proposed for it.

\section{Research on graph-based multi-target scheduling}
\label{sec:tech1-review}

\subsection{The job shop and the flexible job shop}

The classical job shop asks for a sequence of operations on each machine. Job \(j\) has operations \(O_{1j}, O_{2j}, \ldots\) in a fixed order. Operation \(O_{ij}\) has a processing time \(p_{ij}\) and a single eligible machine. Two operations cannot occupy one machine at the same time, and an operation cannot start before its predecessor in the same job has finished. The usual objective is the makespan
\begin{equation}
\cmax = \max_j C_j,
\label{eq:cmax-classic}
\end{equation}
where \(C_j\) is the completion time of the last operation of job \(j\). Garey, Johnson, and Sethi showed that even strongly restricted flow shops and job shops are NP-hard, so a practical factory cannot expect a polynomial algorithm that always returns an optimal sequence \cite{garey1976}. The integer-programming tradition begins with Manne's assignment of binary order variables to pairs of operations that share a machine \cite{manne1960}. Surveys of the deterministic problem run from the shifting-bottleneck procedure of Adams, Balas, and Zawack \cite{adams1988} through the review of Jain and Meeran \cite{jain1999} to the textbook treatment in Pinedo \cite{pinedo2016}.

A school bag factory is not a classical job shop. The body of a backpack can be sewn on either lockstitch machine. A strap can go to the bartack machine or, more slowly, to a lockstitch. Cutting can go to the automatic cutter or to the manual cutter. Brucker and Schlie named this the job shop with multi-purpose machines \cite{brucker1990}. The version studied throughout later work is the flexible job shop scheduling problem, in which each operation \(O_{ij}\) has an eligible set \(M_{ij}\) and a processing time \(p_{ijk}\) on each machine \(k \in M_{ij}\). Brandimarte solved it with tabu search on a routing-and-sequencing representation \cite{brandimarte1993}. Kacem, Hammadi, and Borne made the objective multi-objective, combining makespan with machine workload, and used a localisation heuristic followed by an evolutionary search \cite{kacem2002}. Genetic algorithms with careful encodings became a standard benchmark family \cite{pezzella2008,gao2008,li2016}. Chaudhry and Khan reviewed the technique catalogue and showed how large it had become: dispatching rules, shifting bottlenecks, genetic algorithms, particle swarm methods, and hybrids, almost all of them aimed at makespan on static instances \cite{chaudhry2016}.

The mathematical skeleton that later learning methods still use is the disjunctive graph of Roy and Sussmann \cite{roy1964}. Operations are nodes. A conjunctive arc fixes the route inside a job. A disjunctive pair connects two operations that might share a machine, and a schedule is a choice of one direction for every pair, plus an assignment of flexible operations to machines. Giffler and Thompson had already given a constructive way to build active schedules, in which no operation can be moved earlier without pushing another operation later \cite{giffler1960}. Almost every genetic algorithm and every reinforcement-learning policy in this chapter is, at bottom, a way of choosing those directions.

\begin{figure}[ht]
\centering
\begin{tikzpicture}[
  op/.style={draw, rounded corners, minimum width=1.7cm, minimum height=0.7cm, font=\small},
  mac/.style={draw, fill=gray!12, minimum width=1.7cm, minimum height=0.7cm, font=\small},
  arr/.style={-{Stealth}, thick}
]
\node[op] (a) {cut};
\node[op, right=1.1cm of a] (b) {sew};
\node[op, right=1.1cm of b] (c) {pack};
\node[mac, below=0.9cm of a] (m1) {cutter};
\node[mac, below=0.9cm of b] (m2) {lockstitch};
\node[mac, below=0.9cm of c] (m3) {bench};
\draw[arr] (a) -- (b);
\draw[arr] (b) -- (c);
\draw[arr, dashed] (a) -- (m1);
\draw[arr, dashed] (b) -- (m2);
\draw[arr, dashed] (c) -- (m3);
\end{tikzpicture}
\caption{The classical flexible-job-shop picture. Solid arcs are the route of one school-bag order. Dashed arcs are a machine assignment. Operators, dies, thread lots, colour setups, and machine health are not nodes in this graph.}
\label{fig:classical-graph}
\end{figure}

Figure~\ref{fig:classical-graph} is the picture that most of the papers in this section optimise. It is already enough to make the problem hard. It is not enough for the factory in Section~\ref{sec:problem-set}. The cutter node does not know that a die is required, the lockstitch node does not know which operator is free, and no node carries a wear value. The rest of this section follows the research that made Figure~\ref{fig:classical-graph} smarter, and stops at the point where the missing nodes become the first technology of the thesis.

\subsection{Dispatching rules and why they are still the factory default}

Before metaheuristics and before learning, factories scheduled with priority rules. Panwalkar and Iskander collected more than a hundred of them \cite{panwalkar1977}. Blackstone, Phillips, and Hogg compared the rules that shops actually use \cite{blackstone1982}. Four of them matter for the estimation in Chapter~\ref{ch:estimation}, because they are the baselines a production manager will believe.

First-in-first-out releases work in the order it became ready. It is fair and it is blind. Shortest processing time clears small operations and often keeps the mean flow time down, at the cost of starving a long trolley-bag shell. Earliest due date sorts by the order's promise date and is the natural rule when tardiness is the complaint that reaches the planner. The apparent-tardiness-cost rule of Vepsalainen and Morton combines the two ideas \cite{vepsalainen1987}. For an operation with processing time \(p\) and slack \(s\) it uses a priority of the form
\begin{equation}
Z = \frac{1}{p}\exp\left(-\frac{\max(0,s)}{k\,\bar p}\right),
\label{eq:atc}
\end{equation}
where \(\bar p\) is an average processing time and \(k\) is a look-ahead parameter. A tight, short operation ranks first. A long operation with a comfortable slack waits. Equation~\eqref{eq:atc} is still one of the strongest hand-written rules for weighted tardiness, and it is the strongest of the classical rules one would expect to meet on a due-date-driven bag order book.

Rules have three structural limits that no tuning of \(k\) removes. They score one operation at a time, so they do not see that assigning the navy body to Lockstitch-1 also blocks the only operator who can run the bartack. They do not carry a vector of targets: a weight in front of tardiness and a weight in front of energy are not part of the rule. And they do not look at health. A worn automatic cutter remains the shortest-processing-time choice until it fails, because the rule's \(p\) is the nominal time, not the time at the current wear. The learning methods reviewed next were invented to replace hand-written scores with scores discovered from the shop's own transitions. They inherit the same three limits unless the state is deliberately enlarged.

\subsection{Evolutionary search for several targets}

When a schedule can be encoded as a chromosome, a population can search it. The flexible-job-shop encodings of Pezzella, Morganti, and Ciaschetti \cite{pezzella2008}, Gao, Sun, and Gen \cite{gao2008}, and Li and Gao \cite{li2016} assign each operation a machine and a priority, then decode the chromosome with a list schedule that respects precedence. Hybrids of genetic search and tabu search or variable-neighbourhood descent improve the makespan on the standard Brandimarte and Kacem sets. They are offline procedures. A population of thirty individuals, each decoded by a full shop simulation, is a reasonable night-time calculation and a poor answer to an order that arrives at 10:15 with a same-day promise.

Multi-objective evolutionary algorithms changed the question from ``the best schedule'' to ``the set of schedules that are not beaten on every target.'' Deb, Pratap, Agarwal, and Meyarivan's NSGA-II sorts a population into nondominated fronts and uses crowding distance to keep the front spread out \cite{deb2002}. Zitzler and Thiele's strength Pareto approach, and the later hypervolume indicator, gave the field a way to say whether one front is better than another \cite{zitzler1999,zitzler2003}. When the number of objectives passes three, crowding in NSGA-II stops spreading the population well. Deb and Jain's NSGA-III replaces crowding with reference directions \cite{deb2014}. Gong, De Pessemier, Martens, and Joseph used that algorithm on a flexible job shop whose five objectives were makespan, energy cost, labour cost, peak workload, and total workload, under a time-varying electricity price \cite{gong2019}. Their result is important for this thesis for two reasons. It shows that labour is a real objective, not a soft constraint hidden inside the processing time. And it shows that a many-objective genetic algorithm can expose the conflict between energy and labour, because the cheap electricity hours are often the expensive wage hours.

Energy had already been brought into scheduling by Mouzon, Yildirim, and Twomey, who showed that simply turning idle machines off is not the whole problem: the schedule itself decides when the large loads run \cite{mouzon2007}. Dai, Tang, Giret, and Salido added transportation time and energy to a flexible job shop and solved the bi-objective model with an enhanced genetic algorithm \cite{dai2019}. Sequence-dependent setups, surveyed by Allahverdi and co-authors across two large reviews, sit in the same family \cite{allahverdi2008,allahverdi2015}. A colour change on a sewing machine is exactly such a setup. None of these papers, however, treats the setup, the operator, the energy, and a live wear state as inputs of one online decision.

The practical objection to the evolutionary line, for the factory considered here, is not that the algorithms are weak. On a frozen instance they are often the quality ceiling. The objection is that the ceiling is computed for a shop that holds still. A school-bag week does not hold still. Health moves with every minute of cutting. A rush trolley order is released at mid-morning. A thread lot runs out. Re-evolving a population at every one of these events is possible and usually too slow for the dispatch decision, which is ``what should the free operator start now.'' The learning policies of the next subsection are attractive because, after training, that decision is a forward pass.

\subsection{Reinforcement learning for dispatching}

Reinforcement learning treats scheduling as a sequence of decisions. The state is the shop. The action is the next dispatch. The reward is progress toward the objective. Sutton and Barto's framework is the one used throughout \cite{sutton2018}. Mnih and co-authors showed that a deep network can learn an action value from raw high-dimensional states \cite{mnih2015}. Schulman and co-authors' proximal policy optimisation then became the default stable policy-gradient method for continuous training runs: a clipped ratio keeps a new policy from moving so far that the data just collected become useless \cite{schulman2017}. Haarnoja and co-authors' soft actor-critic adds an entropy bonus so that the policy does not collapse early \cite{haarnoja2018}. Scheduling papers use these algorithms; they do not need to reinvent them. The design problem is the state and the action.

Zhang, Song, Cao, Zhang, Tan, and Xu trained a graph network to dispatch a classical job shop \cite{zhang2020}. The disjunctive graph is the state. A graph neural network embeds every operation. The policy picks the next operation at a machine. Because the network shares weights across nodes, a policy trained on modest instances can be applied to larger ones. That size-agnostic property is the reason later factory work preferred graphs to fixed-length vectors. Park, Chun, Kim, Kim, and Park reported a related graph-and-reinforcement-learning design for job shops \cite{park2021}. Luo studied the flexible job shop when new jobs are inserted, and showed that a deep policy can react without rebuilding a genetic population \cite{luo2020}. Lei, Guo, Zhao, Wang, Qian, Meng, and Tang separated the flexible-job-shop action into several heads, so that operation choice and machine choice are not forced through a single flat list \cite{lei2022}.

These papers moved the field from ``learn a rule for a fixed shop'' to ``learn a rule on a graph that can grow.'' They still optimise, in the typical experiment, one time-based objective, and the graph still contains operations and machines only.

\subsection{Heterogeneous graphs and multi-target policies}

Flexibility makes the graph heterogeneous. An operation is not the same kind of node as a machine, and the relation ``can be processed on'' is not the same relation as ``precedes.'' Song, Chen, Li, and Cao used that fact directly \cite{song2023}. Their IEEE Transactions on Industrial Informatics model extends the disjunctive graph with machine nodes, defines an action as a composite choice of an operation and a machine, and encodes the state with a two-stage heterogeneous graph network. Proximal policy optimisation trains the policy. The published result is that the learned rule beats hand-written rules and still runs quickly on instances larger than the training set. This is the closest ancestor of the first technology in this thesis. It solves the representation problem that plain job-shop graphs do not solve: a machine has its own embedding, and the action is a pair rather than an operation index.

Two limits of that ancestor matter here. The machine node carries scheduling status, such as the time it becomes free, but it does not carry a degradation state that changes the processing time. And there is still no node for the operator or for a shared token. A composite action \((\text{operation},\text{machine})\) is one resource short of a school-bag assignment, which is \((\text{operation},\text{machine},\text{operator})\) and, when the operation needs it, a die or a thread lot.

Wang, Jing, Gu, He, and Chen pushed the same graph idea toward several targets \cite{wang2025}. Their Internet of Things Journal model puts transportation time and energy, together with makespan, into one flexible-job-shop simulation, encodes tasks and machines with a graph network, and trains a multi-agent proximal policy optimisation with a vector-valued critic. The experimental claim is higher solution quality than strong baselines at a small online computing cost. Zeng, Yang, Du, Yang, Wang, and Lv added a dynamic dual-attention network so that a multi-agent policy can react to new jobs and urgent insertions while balancing makespan and average tardiness \cite{zeng2025}. Li, Li, He, and Li reported a collaborative multi-agent method whose reward is split across makespan and energy, with hypervolume used as the comparison metric \cite{li2025}. Attention itself, as a trainable weighted pool of neighbours, comes from the sequence model of Vaswani and co-authors \cite{vaswani2017}, the graph attention network of Velickovic and co-authors \cite{velickovic2018}, and the routing policies of Kool, van Hoof, and Welling \cite{kool2019}. Graph convolutional networks \cite{kipf2017} and GraphSAGE \cite{hamilton2017} are the earlier message-passing tools that these attention schedulers refine.

Table~\ref{tab:sched-lineage} places the lineage against the decision a bag factory actually needs.

\begin{table}[ht]
\centering
\caption{What recent scheduling models decide, and what they leave out of the school-bag assignment.}
\label{tab:sched-lineage}
\small
\begin{tabularx}{\textwidth}{@{}lXXXX@{}}
\toprule
Work & Decision & Targets & Live health & Shared tokens \\
\midrule
Brandimarte \cite{brandimarte1993} & route and sequence & makespan & no & no \\
Kacem et al. \cite{kacem2002} & route and sequence & time and load & no & no \\
Gong et al. \cite{gong2019} & route, sequence, labour & five, including energy and labour & no & labour hours \\
Zhang et al. \cite{zhang2020} & next operation & makespan & no & no \\
Song et al. \cite{song2023} & operation and machine & makespan & no & no \\
Wang et al. \cite{wang2025} & operation and machine & makespan, transport, energy & no & no \\
Zeng et al. \cite{zeng2025} & two agents & makespan and tardiness & no & no \\
This thesis & operation, machine, operator, token & time, due date, energy, balance, failures & yes, through Chapter~\ref{ch:methods}'s twin & dies and thread lots \\
\bottomrule
\end{tabularx}
\end{table}

The last row is not a claim that earlier rows are incomplete as papers. Each of them answers the question it asked. The row says that the question asked by a school-bag planner is wider than any one of those questions, and that the widening is specific: the action must name a resource bundle, and one of the node features must be allowed to come from a machine-status model. That is the design brief of resource-coupled heterogeneous graph attention scheduling, written out in Section~\ref{sec:rch}.

\subsection{Textile shops and the second resource}

General flexible-job-shop benchmarks are machining benchmarks. A school bag is closer to a sewing shop. Chen, Xie, Ma, Chen, and Wang formulated a textile flexible job shop and modified ant-colony optimisation so that the pheromone remembers the product as well as the machine, with a heuristic that lifts machines which the search would otherwise ignore \cite{chen2024}. The objective is balance of machine time across the operations of the textile route. The paper is useful because it shows that a textile route is not an arbitrary chain: machines and product types constrain each other, and a search that ignores the product will over-use the convenient machine. It does not model operators as competing resources, and it does not model health.

Sewing research has a name for the operator constraint. A dual-resource-constrained shop is one in which a station needs both a machine and a worker, and workers are fewer or less flexible than machines. Xu, Xu, and Xie reviewed that literature \cite{xu2011}. The consequence for dispatching is immediate. A machine that is free in the data table may be unusable because the only qualified operator is inside another operation. Sequence-dependent setups apply to the operator as well as to the machine: a change of colour, of needle, or of bag style costs minutes of the person, not only of the head. Gong and co-authors' labour-aware model \cite{gong2019} prices this effect. It still plans labour as a cost accumulated by a static chromosome, not as a token the online policy must seize and release.

Rework belongs in the same practical family. Flapper, Fransoo, Broekmeulen, and Inderfurth reviewed planning when a quality failure sends work back \cite{flapper2002}. A school-bag inspection can fail a zipper or a crooked strap. The estimation shop in Chapter~\ref{ch:estimation} keeps inspection as a required operation rather than a random rework loop, so that the comparison of schedulers is not mixed with a second random process. The formulation in Chapter~\ref{ch:methods} leaves a place for a rework arc; the numerical shop does not draw one.

\subsection{The gap that the first technology is built to close}

The scheduling literature, read against a school-bag route, supports five conclusions.

The problem is a flexible job shop, and it is reasonable to keep the disjunctive-graph view of precedence and machine choice \cite{roy1964,song2023}. The action, however, is larger than a machine assignment. It includes an operator whose skill changes the processing time, and a token whose capacity is one or two, not the number of machines. A graph that has no operator node and no token node cannot attend to the contention that actually blocks the hardware bench and the khaki lot.

The targets are plural. Makespan research \cite{brandimarte1993,pezzella2008,zhang2020,song2023} is the right foundation and the wrong single number. Energy \cite{mouzon2007,dai2019,gong2019,wang2025}, due dates \cite{vepsalainen1987,zeng2025}, and workload balance \cite{kacem2002,chen2024,gong2019} have each been studied. A factory policy has to carry them together, and a manager has to be able to move the emphasis without training an entirely new system. Preference-conditioned decisions are the natural shape of that requirement. The recent multi-agent critics \cite{wang2025,li2025,zeng2025} move in this direction and still do not condition a resource-bundle action on a preference vector tied to bag-factory tokens.

The useful online object is a learned dispatch policy, not a population that is re-evolved at every breakdown. Graph policies are the ones that transfer across order-book sizes \cite{zhang2020,song2023}. They are also the ones that can accept an extra node type without a new hand-written feature hierarchy.

Setups are sequence dependent and, in this factory, mostly colour changes \cite{allahverdi2015}. A policy that does not see the previous colour on the machine will pay the setup in the simulation and never in the decision.

No paper in this section writes the machine's wear into the same graph that chooses the operator and the thread lot. That omission is not a small missing feature. It is a second technology, reviewed next, which this thesis couples to the first instead of leaving it as a report beside the plan.

\section{Research on machine status and digital twins}
\label{sec:tech2-review}

\subsection{Status is a signal, not a label}

A machine status that a scheduler can use is a number with a clock. ``The cutter is fine'' is not such a number. The cutter's vibration, the current of its feed motor, the temperature at the knife, the count of thread breaks on a sewing head, and the hours since the last knife change are such numbers. They arrive from a controller, they are noisy, and they are late by whatever period the collector uses.

Two standards make this more than a vendor choice. IEC 62264, the standard built on ISA-95, separates the enterprise plan from the control system and defines the objects that pass between them: equipment, material lots, personnel, and process segments \cite{iec62264}. IEC 62541, OPC Unified Architecture, is the usual way those equipment objects expose live values to a client outside the programmable controller \cite{iec62541}. A scheduler that claims to react to machine status is, in factory terms, a client of this stack. It does not need to own the sensors. It needs a defined place where a health value enters the decision.

Lee, Bagheri, and Kao organised the same idea as a five-level architecture for cyber-physical manufacturing: a smart connection, conversion of data to information, a cyber model, cognition, and configuration of the physical system \cite{lee2015}. Monostori and co-authors reviewed cyber-physical production systems as the place where this loop has to be closed, and listed scheduling among the decisions that still sit too far from the physical state \cite{monostori2016}. Xu, Xu, and Li's survey of Industry 4.0 makes the same point at the level of the whole plant: connectivity is no longer the scarce piece, and decision models that use the connectivity are \cite{xu2018}. Koren and Shpitalni's reconfigurable manufacturing systems add a design consequence. If the plant can change its routing, the schedule is part of the configuration, not a file written once per shift \cite{koren2010}.

The school-bag factory does not need a new network architecture. It needs the cognition level of Lee's stack to be a scheduler, and the cyber level to be a degradation model whose output is a feature of that scheduler. The next subsections review the models that can sit at that cyber level.

\subsection{Digital twins of a production system}

A digital twin is a living model paired with a particular physical system, updated by that system's data, and used to make a decision about that system. Grieves and Vickers stated the idea as a way to see undesirable behaviour before the physical system has to exhibit it \cite{grieves2017}. Tao, Zhang, Liu, and Nee reviewed industrial practice and separated a digital twin from a plain simulation by the requirement of ongoing synchronisation \cite{tao2019}. Qi and Tao compared the twin with big-data manufacturing and argued that the twin is the structure that makes data actionable \cite{qi2018}. Negri, Fumagalli, and Macchi reviewed the role of twins inside cyber-physical production and found a wide spread of meanings, from geometric models to scheduling models \cite{negri2017}. Cimino, Negri, and Fumagalli followed the manufacturing applications and found the same spread \cite{cimino2019}. Fuller, Fan, Day, and Barlow listed enabling technologies and the open problems, among them the difficulty of keeping the model and the asset in step \cite{fuller2020}. Lu, Liu, Wang, Huang, and Xu proposed a reference model for a twin-driven manufacturing system and again placed production planning among the applications that the reference model must be able to host \cite{lu2020}.

For this thesis the word ``twin'' is used in a narrow way. The twin of a school-bag machine is not a three-dimensional drawing of the cutter. It is a state \(h_k(t)\) for each machine \(k\), a rule that updates \(h_k\) from the latest signals, and a pair of maps from \(h_k\) to the quantities the scheduler needs: a processing-time multiplier and a probability that the next cycle is interrupted. If those maps are not in the scheduling decision, the factory has a dashboard, not a twin of the plan.

\begin{figure}[ht]
\centering
\begin{tikzpicture}[
  box/.style={draw, rounded corners, align=center, minimum width=2.6cm, minimum height=0.95cm, font=\small},
  arr/.style={-{Stealth}, thick}
]
\node[box] (phys) {Physical\\machine};
\node[box, right=1.3cm of phys] (sig) {Signals\\OPC~UA};
\node[box, right=1.3cm of sig] (twin) {Health\\state \(h_k\)};
\node[box, right=1.3cm of twin] (plan) {Schedule\\decision};
\draw[arr] (phys) -- (sig);
\draw[arr] (sig) -- (twin);
\draw[arr] (twin) -- (plan);
\draw[arr] (plan.south) -- ++(0,-0.7) -| (phys.south);
\end{tikzpicture}
\caption{The twin loop required by a status-aware schedule. The return arrow is the part most digital-twin case studies leave as a manual action by the planner.}
\label{fig:twin-loop}
\end{figure}

Figure~\ref{fig:twin-loop} is deliberately small. The research question is the return arrow. Many twin deployments stop at the health state and show it to a person. The person then edits the plan. That can be the right engineering choice for a rare failure. It is a poor choice for a wear process that moves a little on every backpack, because the plan is then permanently out of date by the amount the person has not yet typed.

\subsection{Degradation and remaining useful life}

Remaining useful life is the time left until a machine crosses a failure threshold. Si, Wang, Hu, and Zhou reviewed the statistical approaches and separated models that regress a health indicator from models that write an explicit stochastic process for the degradation \cite{si2011}. Lei, Li, Guo, Li, Yan, and Lin reviewed the path from sensors to a life prediction, including the feature extraction that a factory must get right before any scheduler sees a number \cite{lei2018}. This thesis does not compete with that extraction literature. It assumes a scalar health indicator has already been formed, and it spends its effort on the use of the scalar.

A Wiener process is the usual explicit model when degradation can be treated as a drift plus noise. In the linear form,
\begin{equation}
h(t) = h(0) + \eta t + \sigma B_t,
\label{eq:wiener-review}
\end{equation}
where \(\eta\) is the wear rate, \(\sigma\) is a diffusion scale, and \(B_t\) is standard Brownian motion. The increment over a cycle of length \(\Delta\) is Gaussian with mean \(\eta\Delta\) and variance \(\sigma^2\Delta\). Si, Wang, Hu, Zhou, and Pecht extended the idea to nonlinear drift, which matters when wear accelerates near failure \cite{si2012}. Gebraeel, Lawley, Li, and Ryan put a Bayesian update on the residual-life distribution so that each new degradation signal revises the life left, rather than only the point estimate \cite{gebraeel2005}. Kalman filtering is the classical linear-Gaussian way to carry such a revision \cite{kalman1960}. Particle filters cover the non-Gaussian case at a higher computational cost \cite{arulampalam2002}.

The scheduling consequence of Equation~\eqref{eq:wiener-review} is a chance constraint, not a single predicted day of failure. Charnes and Cooper introduced chance constraints as a way to keep a random inequality feasible with a chosen probability \cite{charnes1959}. Birge and Louveaux place that idea inside stochastic programming \cite{birge2011}. For a machine, the natural statement is
\begin{equation}
\Prob\big(h_k(t+\Delta) \ge h^{\star}\big) \le \varepsilon,
\label{eq:chance-review}
\end{equation}
with \(\varepsilon\) chosen by the factory. If the increment is Gaussian, Equation~\eqref{eq:chance-review} becomes a deterministic threshold on the current health, the cycle length, and \(\varepsilon\). Chapter~\ref{ch:methods} writes that threshold down. The point for the present review is that the constraint is old, the degradation model is mature, and their appearance inside a graph-attention dispatcher is not.

Nguyen and Medjaher connected a deep prognostic model to a dynamic predictive-maintenance decision, so that the maintenance action, rather than only the life number, is the output \cite{nguyen2019}. Wang, Li, and Zhao built multi-objective maintenance schedules that take a remaining-useful-life prediction as an input and trade maintenance duration against cost \cite{wang2024cie}. Both papers treat maintenance as the decision. A production schedule, if present, is outside the model or is a fixed demand the maintenance plan must serve.

\subsection{Maintenance inside the production schedule}

The closer line of work inserts maintenance into the job shop itself. A 2024 study of bi-objective flexible job shop scheduling with condition-based maintenance updates the degradation parameters by Bayesian inference after each maintenance, and uses an evolutionary algorithm with an insertion step so that a maintenance window disturbs the production plan less than a naive split would \cite{cbm2024}. The objectives are production and maintenance together. The algorithm is offline. Machine health changes the feasibility of a chromosome; it is not a feature of a learned online policy, and operators are not part of the chromosome.

Pandhare, Negri, Ragazzini, Cattaneo, Macchi, and Lee described a prognostics-enabled twin that schedules a flow shop while equipment is degrading \cite{pandhare2024}. Field data, reduced by principal components, drive a discrete-event model. A genetic algorithm searches the schedule. The reported comparison is that a twin which knows the degraded state estimates makespan more realistically than a twin which pretends the equipment is nominal, and that the resulting schedule is better for that reason. The paper is the clearest statement that health belongs inside scheduling. It is a flow shop, so routing flexibility is absent. It has no operator tokens and no thread lots. The search is a genetic algorithm, so a rush order or a sudden jump in a health residual means another population, not a masked action in a policy that is already trained.

An online joint scheme for predictive maintenance and production planning, of the kind discussed by Zhu, Wang, and Yuan, uses a rolling horizon: at each period the latest remaining-useful-life estimate and the latest demand forecast are written into a shorter optimisation, and only the first period is implemented \cite{zhu2024pdm}. Rolling horizons are the right operational pattern for a planner who re-solves. They are heavy if the inner model is a full flexible job shop with resource tokens, and they are unnecessary if a trained policy can be evaluated at the epoch when the twin says the residual has moved enough.

The maintenance literature and the scheduling literature therefore meet, but they meet in evolutionary or rolling-horizon optimisers. They do not meet in the heterogeneous graph policy of Song, Chen, Li, and Cao \cite{song2023} or in the multi-target policy of Wang, Jing, Gu, He, and Chen \cite{wang2025}. The second technology of this thesis is that meeting. It is specified in Section~\ref{sec:ssdt} as a status-synchronized twin whose health, failure probability, and maintenance action are features and actions of the same dispatcher.

\subsection{The gap that the second technology is built to close}

Four observations close the review.

Health is a state variable with a stochastic increment, not a static penalty. Wiener and Bayesian residual-life models already supply the increment \cite{si2011,si2012,gebraeel2005}. A scheduler that reads a health value once per shift is ignoring the model it paid to build.

The map from health to the shop has two outputs, and both belong in the decision. The processing time grows as the machine leaves its nominal condition, and the probability of an interruption grows with it. Energy grows as well, because a worn drive draws more for the same piece. A makespan model with nominal times will prefer the fast worn cutter and then spend the saved minutes on a repair.

The maintenance action has to be cheap enough, and visible enough, to be chosen by the same policy that chooses the backpack. A separate maintenance optimiser that runs overnight will insert windows that the morning rush has already invalidated. A chance constraint of the form \eqref{eq:chance-review}, evaluated at the dispatch epoch, is the lightest coupling that still refuses a dangerous assignment \cite{charnes1959}.

The twin must be allowed to force a new decision. Synchronisation that only refreshes a dashboard does not change \(C_j\). Event-triggered rescheduling, called when the predicted failure probability or a newly released order crosses a threshold, is the configuration step in the cyber-physical stack of Lee, Bagheri, and Kao \cite{lee2015}. The trigger condition is part of the technology, not an implementation detail.

\section{Problem set}
\label{sec:problem-set}

\subsection{The school-bag route}

The product family in this thesis is the ordinary school bag: a soft-sided pack carried by a student, together with two cousins that share the same shop, a trolley bag and a sports sack. The factory is imagined as a single-site cell, not a supply chain. Cloth, webbing, zippers, and hardware arrive as lots. Finished bags leave in cartons. The planning horizon of the estimation is one congested day, long enough for wear and colour changes to matter and short enough that the order book is a list rather than a season.

A classic student backpack crosses eight steps. The main shell is cut. The lining is cut. A logo is embroidered or printed. The body, including the compartment seams, is sewn. The shoulder straps and handle are sewn, preferably on a bartack or heavy lockstitch because the webbing is thick. The zipper, buckles, and adjusters are attached at the hardware bench. The bag is inspected. It is packed. The two cutting steps both need a die. Embroidery and the sewing steps need a thread lot of the order's colour. Hardware, inspection, and packing do not.

A trolley bag crosses seven steps: cutting the shell, reinforcing the base, sewing the shell, fitting the frame, fitting the wheels, inspection, and packing. The frame and the wheels share the hardware bench, so a trolley order holds that bench longer than a backpack does. A sports sack is the short route: cut, print, sew, drawstring, inspect, and pack. It is the order a shortest-processing-time rule will always prefer, and the order that fills the embroidery head while a trolley shell waits.

The eligible machines are deliberately uneven. Two cutters share the cutting group. The automatic cutter is fast, high-powered, and, in the case studied in Chapter~\ref{ch:estimation}, already worn at the start of the day. The manual cutter is slow, cheap in power, and healthy. One embroidery head serves every logo. Two lockstitch machines and one bartack machine serve the sewing and reinforcement steps, with a strap allowed to fall back from the bartack onto a lockstitch. One hardware bench, one inspection table, and one packing table finish the route. This mix is a partial flexible job shop in the sense of Brandimarte \cite{brandimarte1993}: some operations have a real choice, and some have only one machine.

\begin{figure}[ht]
\centering
\begin{tikzpicture}[
  st/.style={draw, rounded corners, align=center, font=\scriptsize, minimum height=0.72cm, minimum width=1.15cm},
  arr/.style={-{Stealth}, thin}
]
\node[st] (c1) {cut\\shell};
\node[st, right=0.35cm of c1] (c2) {cut\\lining};
\node[st, right=0.35cm of c2] (pr) {logo};
\node[st, right=0.35cm of pr] (sb) {sew\\body};
\node[st, right=0.35cm of sb] (ss) {sew\\strap};
\node[st, right=0.35cm of ss] (hw) {hardware};
\node[st, right=0.35cm of hw] (qc) {inspect};
\node[st, right=0.35cm of qc] (pk) {pack};
\draw[arr] (c1) -- (c2);
\draw[arr] (c2) -- (pr);
\draw[arr] (pr) -- (sb);
\draw[arr] (sb) -- (ss);
\draw[arr] (ss) -- (hw);
\draw[arr] (hw) -- (qc);
\draw[arr] (qc) -- (pk);
\end{tikzpicture}
\caption{Route of one classic school backpack. Trolley bags and sports sacks use shorter or longer chains on the same cells. Cutting and colour-sensitive sewing are the steps that seize a shared token.}
\label{fig:backpack-route}
\end{figure}

Colour is a setup, not only a bill-of-material attribute. If a sewing machine or the embroidery head moves from navy to red, the operator changes thread and the machine loses a setup of about eleven minutes in the case data. A repeated colour pays a short warm-up only when the machine has been idle and cleared. Allahverdi's setup surveys describe exactly this structure \cite{allahverdi2015}. In a bag factory the setup is large relative to a small sports sack and small relative to a heavy trolley shell, so a policy that groups colour is not automatically a policy that protects the trolley due date. That conflict is one reason a single rule is not enough.

Quantity and size scale the cycle. The nominal times in Chapter~\ref{ch:estimation} are written for a batch of eighty bags of size M. A batch of one hundred and twenty, or a size L pattern, takes longer on every step and occupies the token for longer. The die and the thread lot are therefore not only feasibility constraints. They are congestion multipliers.

\subsection{Resources and the information the scheduler is allowed to see}

The shop is dual-resource constrained \cite{xu2011}. Eight operators are modelled in the estimation cell. Two can cut, and the second of them is slower. One runs embroidery and can help at inspection. Three sew; only two of those three are qualified on the bartack, and their skills differ. One operator owns the hardware bench and can pack. One owns inspection and can also pack or, slowly, work the hardware bench. An operation is feasible only if at least one qualified operator is free at the same moment as an eligible machine. Skill scales the processing time. A weak pairing is legal and expensive.

Shared tokens sit on top of the operator constraint. Two cutting dies can be out at once, so the two cutters can run together only when two dies are free. Thread lots are counted per colour: two concurrent users for navy, black, red, and grey, and one concurrent user for khaki. A shortage scenario in Chapter~\ref{ch:estimation} drops navy, red, and khaki to a single lot. The token is seized when the operation starts and released when it ends. A scheduler that ignores the token will recommend two khaki bodies at once and then discover that the second cannot start.

The information model is intentionally modest. At a decision epoch the scheduler may see the released orders, their routes, colours, sizes, quantities, and due dates; the eligibility and skill tables; the previous colour and the free time of each machine and operator; the token counts; and a scalar health for each machine. It may not see the future rush order before its release time. It may not see a breakdown that has not yet been implied by the health model. This is the information a manufacturing-execution client can actually subscribe to under IEC 62264 and OPC UA \cite{iec62264,iec62541}. It is also enough, as Chapter~\ref{ch:methods} shows, to build the graph.

\subsection{Targets}

The planner does not have one number to minimise. Five quantities describe a finished schedule \(S\).

The makespan \(\cmax(S)\) is the clock time at which the last bag of the released book is packed. It measures the length of the day.

The total tardiness
\begin{equation}
T(S) = \sum_j \max(0, C_j - d_j)
\label{eq:tard-problem}
\end{equation}
measures broken promises. A sports sack that finishes an hour early does not cancel a trolley bag that finishes an hour late. On-time rate, the fraction of orders with \(C_j \le d_j\), is reported beside \(T\) because a single very late order can dominate the sum.

Energy \(E(S)\) is the electrical energy of processing, setup, preventive maintenance, and the expected energy of a repair. Power is a property of the machine. A worn machine draws more for the same minute. The automatic cutter dominates this sum when it is used carelessly.

Workload imbalance \(B(S)\) is the standard deviation of busy time across machines. A schedule that parks almost all cutting on one cutter will look acceptable on makespan until that cutter becomes the wear problem of the next day. Balance is therefore both a fairness target and a leading indicator of tomorrow's health.

Expected failures \(\Phi(S)\) sum, over every executed cycle, the probability that the cycle is interrupted. The estimation uses this expectation rather than a single random breakdown draw, so that two schedulers are compared on the risk they chose and not on a lucky or unlucky random seed. Expected repair delay is accumulated into the timeline for the same reason.

These five numbers conflict. The fast cutter reduces the nominal cutting time and increases energy, wear, and failure probability. Grouping navy orders reduces setups and can increase the tardiness of a red order. Preventive maintenance spends sixteen minutes now to avoid an expected repair of several times that, and to restore a faster, cooler cycle. A preference vector \(\omega\), with non-negative components summing to one, tells the policy which of the conflicts matters for the present shift. Changing \(\omega\) is a management action. It should not require a new research project.

\subsection{Why the previous methods do not close this problem}

Put an earliest-due-date rule on this shop and it will protect promises better than first-in-first-out, as Chapter~\ref{ch:estimation} confirms. It will also send the logo of every urgent order to the single embroidery head without seeing the queue it is building, and it will keep the worn automatic cutter busy because the rule has no health input. The apparent-tardiness-cost rule improves the local ranking and still has no token, no skill, and no wear.

Put NSGA-II or NSGA-III on a frozen snapshot and the front can be excellent by morning \cite{deb2002,deb2014,gong2019}. At 10:15 a rush order and a jumped health residual make the snapshot false. Rebuilding the front is a batch computation. The dispatch decision is not.

Put the heterogeneous graph policy of Song, Chen, Li, and Cao on the shop and the operation-machine choice becomes learned and fast \cite{song2023}. The operator may still be double-booked by a later repair of the action, because the operator was never in the graph. The khaki lot may be double-booked for the same reason. Health may be displayed on a screen in the same room and never enter the embedding.

Put a condition-based maintenance optimiser beside the plan \cite{cbm2024,pandhare2024,wang2024cie} and the cutter will be maintained on a sensible wear threshold. The maintenance window will not know that the operator it needs is the only person who can start the urgent strap, and the production policy will not know, in its features, that the maintenance has just made the cutter fast again.

The problem set of this thesis is the joint problem these methods each solve half of. Find a dispatch policy that, at every epoch, selects a feasible bundle \((\text{operation},\text{machine},\text{operator},\text{token})\) so that a preference-weighted combination of makespan, tardiness, energy, imbalance, and expected failures is small, with processing times and failure probabilities written by a health state that is updated from machine status, and with a maintenance action available in the same decision when a chance constraint would otherwise be broken.

\subsection{The two technologies, stated as contributions}

The first technology is resource-coupled heterogeneous graph attention scheduling (RCH-GAT). The shop graph contains four node types that the assignment actually touches: the operation, the machine, the operator, and the shop-level contention for tokens. Relation-specific attention pools those nodes into a score. The score is conditioned on the preference vector \(\omega\). Infeasible bundles, including bundles that fail skill, precedence, machine capacity, or token capacity, are masked before the score is used. The training objective in the full model is proximal policy optimisation of this masked, preference-conditioned policy. The estimation system of Chapter~\ref{ch:estimation} implements the same score in a one-hop form whose coefficients are identified by cross-entropy search, so that every reported number comes from an executable policy rather than from an unreported training run. Section~\ref{sec:rch} writes the full model. The identification procedure is part of the system build.

The second technology is the status-synchronized digital twin with chance-constrained rescheduling (SSDT-CC). Each machine carries a health state. The state scales processing time and energy, and it supplies the failure probability of the cycle about to be started. If that probability exceeds a tolerance \(\varepsilon\), or if the predicted health at the end of the cycle exceeds a wear ceiling, the decision inserts a short preventive maintenance, renews the health, and recomputes the cycle. The policy is an online policy, so a new health value is visible at the next dispatch without a separate rescheduling project. The chance constraint is what turns the twin from a display into a veto. Section~\ref{sec:ssdt} writes the state, the constraint, and the coupling back into the graph features of the first technology.

The two technologies are one system. RCH-GAT without the twin can still prefer a healthy machine, but it cannot renew a worn one. The twin without RCH-GAT can still maintain the cutter, but it will maintain whatever machine a shortest-processing-time rule happens to overload. Chapter~\ref{ch:estimation} estimates them separately and together for that reason.

\subsection{Questions the estimation is required to answer}

The estimation is not a collection of illustrations. It is required to answer four questions about the problem just stated.

First, on a nominal school-bag book, does the coupled system reduce makespan, energy, imbalance, and expected failures relative to first-in-first-out, shortest processing time, earliest due date, and apparent tardiness cost, without giving those rules a large advantage on tardiness? A system that saves energy by missing every due date has not solved a multi-target problem.

Second, when the day is stressed, does the same system hold up? The stresses used in Chapter~\ref{ch:estimation} are a mid-day rush release, an accelerated wear regime, and a thread-lot shortage. These are the three ways the problem set says the shop escapes a static flexible-job-shop benchmark.

Third, which technology produces which part of the gain? Removing resource features and removing status-and-maintenance are the two ablations. If both removals do nothing, the technologies are not doing the work the model description claims.

Fourth, does the chance-constraint tolerance behave like a real knob? Tightening \(\varepsilon\) should buy fewer expected failures at the price of more preventive maintenance. If the estimation cannot show that exchange, the constraint is only a story.

A fifth, softer question is how stable the tardiness comparison is. Total tardiness is a sum of a few late orders, so its standard deviation across random books is expected to be large. The estimation reports that spread instead of treating a small difference in the mean as a ranking.

\subsection{How the remaining chapters answer the problem set}

Chapter~\ref{ch:methods} answers the problem set with formulas. Section~\ref{sec:rch} writes objectives, constraints, the heterogeneous graph, the attention score, the masked action, and the policy-gradient training of RCH-GAT, and it walks one numerical assignment. Section~\ref{sec:ssdt} writes the health dynamics, the Kalman revision of that health, the chance constraint and its Gaussian threshold, the maintenance map, and the way the renewed health re-enters the graph.

Chapter~\ref{ch:estimation} answers the problem set with a system and a measurement. Section~\ref{sec:build} describes the shop that was built, the data it accepts, and the executable form of the two technologies. Section~\ref{sec:estimate} reports the four scenarios, the ablations, the tolerance sweep, and the limits of what those runs can support. The program that produced the tables is the program in the thesis directory. The numbers were not taken from an external factory and they were not taken from the papers reviewed in this chapter.

\chapter{The Two Technologies}
\label{ch:methods}

The problem set in Section~\ref{sec:problem-set} asks for one system and names two technologies inside it. This chapter specifies both of them at the level a second implementation could follow. Section~\ref{sec:rch} specifies resource-coupled heterogeneous graph attention scheduling: the objectives, the resource constraints, the graph, the attention score, the masked policy, and the way that policy is trained. Section~\ref{sec:ssdt} specifies the status-synchronized digital twin: the health dynamics, the revision of health from a noisy reading, the chance constraint, the preventive-maintenance action, and the write-back of the renewed health into the graph of Section~\ref{sec:rch}. Each section ends with the same numerical assignment, a navy backpack shell on the two cutters, so that the formulas are not left as symbols only.

The estimation in Chapter~\ref{ch:estimation} uses a one-hop executable form of the score and a hazard form of the chance constraint. Those forms are stated here as the operational model. The multi-layer network and the Kalman revision are stated as the plant model one would run when real sensor streams, rather than a simulator that already knows its own wear, are the source of \(h_k\).

\section{Resource-coupled heterogeneous graph attention}
\label{sec:rch}

\subsection{Sets, routes, and the decision that is actually taken}

Let \(J\) be the set of released orders and \(M\) the set of machines. Each order \(j \in J\) has a release time \(r_j\), a due date \(d_j\), a colour \(c_j\), a size, a quantity \(q_j\), and an ordered route of operations \(O_{1j},\ldots,O_{n_j j}\). Operation \(O_{ij}\) has a base duration \(\tau_{ij}\) defined at the reference quantity \(q_0 = 80\) and reference size M, an eligible set of machine groups \(G_{ij}\), and a token type \(u_{ij} \in \{\mathrm{die},\mathrm{colour},\mathrm{none}\}\).

Let \(K\) be the set of operators. Operator \(o\) has a skill \(\sigma_{ok} > 0\) only for the machine groups that person is allowed to run, and \(\sigma_{ok} = 0\) otherwise. Let \(Q\) be the set of token pools. The die pool has capacity \(n_{\mathrm{die}}\). The colour pool of colour \(c\) has capacity \(n_c\). A pool of capacity zero would make every operation of that colour infeasible; the factory data do not contain such a pool.

A dispatch at time \(t\) is a bundle
\begin{equation}
a = (i,j,k,o) ,
\label{eq:bundle}
\end{equation}
meaning that operation \(O_{ij}\) starts on machine \(k\) with operator \(o\). If \(u_{ij}\) is a die or a colour, the bundle also seizes one unit of that pool for the whole occupation of the machine. The bundle is the action of the policy. It is larger than the operation-machine pair of Song, Chen, Li, and Cao \cite{song2023}, and the extra components are the reason the graph in this section has more than two node types.

\subsection{Objectives}

A complete schedule \(S\) is the set of bundles that were started, together with the clock times at which they started and finished. Five functions of \(S\) are kept separate until a preference vector combines them. The makespan is
\begin{equation}
f_1(S) = \cmax = \max_{j \in J} C_j .
\label{eq:f1}
\end{equation}
Total tardiness is
\begin{equation}
f_2(S) = \sum_{j \in J} \max(0, C_j - d_j) .
\label{eq:f2}
\end{equation}
Energy sums every occupation, including setup minutes and preventive-maintenance minutes, and adds the expected energy of an interruption,
\begin{equation}
f_3(S) = \sum_{a \in S} P_k \frac{\Delta_a}{60}\big(1 + \gamma_E h_k\big) + \sum_{a \in S} \phi_a \, P_k^{\mathrm{rep}} \frac{\Delta^{\mathrm{rep}}}{60}\big(1 + \gamma_E h_k\big) .
\label{eq:f3}
\end{equation}
Here \(P_k\) is the nameplate power of machine \(k\) in kilowatts, \(\Delta_a\) is the occupation in minutes, \(h_k\) is the health used for that occupation, \(\gamma_E = 0.45\) is the wear surcharge on power, \(\phi_a\) is the failure probability of the cycle, and the repair term uses half the nameplate power for \(\Delta^{\mathrm{rep}} = 48\) minutes. The half-power repair is a modelling choice: a stopped head does not cut, but the drive and the recovery move are not free.

Workload imbalance is the standard deviation of busy minutes,
\begin{equation}
f_4(S) = \left(\frac{1}{|M|}\sum_{k \in M}\big(B_k - \bar B\big)^2\right)^{1/2}, \qquad B_k = \sum_{a:\,k(a)=k} \Delta_a .
\label{eq:f4}
\end{equation}
Expected failures are the sum of the cycle probabilities,
\begin{equation}
f_5(S) = \sum_{a \in S} \phi_a .
\label{eq:f5}
\end{equation}
A preference \(\omega \in \R^4\) with \(\omega \ge 0\) and \(\sum_{m=1}^{4}\omega_m = 1\) defines the scalar that a single policy optimises,
\begin{equation}
J_\omega(S) = \sum_{m=1}^{4} \omega_m \, \tilde f_m(S) + \lambda_\Phi \, \tilde f_5(S) ,
\label{eq:scalar}
\end{equation}
where a tilde denotes division by the same function evaluated on the first-in-first-out schedule of the same order book. The ratio keeps a minute of makespan from being added to a kilowatt-hour. The failure term keeps a weight \(\lambda_\Phi\) even when the manager's four-way preference does not mention breakdowns, because a schedule that fails machines is not an admissible way to win the other four numbers. In the estimation, \(\lambda_\Phi = 0.12\) during identification and the reported tables always show the five raw functions, not only \(J_\omega\).

Equation~\eqref{eq:scalar} is a scalarisation, not a claim that the five functions have a single optimum. Changing \(\omega\) is how the factory moves along the conflicts described in Section~\ref{sec:problem-set}. NSGA-II remains available as an offline reference precisely because a scalarisation can miss a region of the front \cite{deb2002}. The online system is allowed to be a scalarisation because the dispatch decision cannot wait for a front.

\subsection{Constraints}

Let \(x_{ijko} \in \{0,1\}\) be one when bundle \((i,j,k,o)\) is the bundle that executes \(O_{ij}\). Every operation is executed once,
\begin{equation}
\sum_{k \in M_{ij}} \sum_{o \in K_{ij}} x_{ijko} = 1 ,
\label{eq:assign-once}
\end{equation}
where \(M_{ij} = \{k : \mathrm{group}(k) \in G_{ij}\}\) and \(K_{ij}\) is the set of operators with positive skill on the chosen machine. Precedence inside the route is
\begin{equation}
S_{i+1,j} \ge C_{ij} , \qquad C_{ij} = S_{ij} + \Delta_{ij} ,
\label{eq:precedence}
\end{equation}
and the release date is \(S_{1j} \ge r_j\).

Machine capacity and operator capacity are disjunctive. For any two distinct occupations of the same machine, one finishes before the other starts, and the same statement holds for an operator. With a binary order variable \(y\),
\begin{equation}
C_{a} \le S_{a'} + H(1-y_{aa'}) , \qquad C_{a'} \le S_{a} + H y_{aa'} ,
\label{eq:disjunct}
\end{equation}
for a horizon constant \(H\) large enough to dominate any feasible day. Token capacity is a knapsack at every time. For the die pool, and likewise for each colour with its own capacity,
\begin{equation}
\sum_{a \in A(t)} \1\{u_a = \mathrm{die}\} \le n_{\mathrm{die}} ,
\label{eq:token}
\end{equation}
where \(A(t)\) is the set of occupations that cover time \(t\). Skill is a hard mask rather than a penalty: \(x_{ijko} = 0\) whenever \(\sigma_{o,\mathrm{group}(k)} = 0\).

Setups are sequence dependent on colour. If machine \(k\) has just run colour \(c\) and the next occupation has colour \(c_j \ne c\), a setup of \(s^{\mathrm{col}} = 11\) minutes is inserted before processing. The first occupation of a colour-sensitive machine after a cleared state pays a warm-up of \(s^{\mathrm{warm}} = 4\) minutes. Hardware, inspection, and packing do not pay a colour setup. The setup occupies both the machine and the operator, and it occupies the token, because the die or the thread is the thing being changed.

These constraints are the integer program a solver would be given if the health state were frozen and no new order could appear. The shop does not offer those two favours. Health is a function of the schedule, so \(\Delta_{ij}\) depends on a state that the binaries themselves change, and orders appear at release dates the morning plan does not know. The integer program is the definition of feasibility. The policy below is the object that constructs a feasible \(x\) online.

\subsection{Processing time at a given health}

The duration of the processing part of a bundle, before setup and before any maintenance, is
\begin{equation}
p_{ijko}(h_k) = \tau_{ij}\,\frac{q_j}{q_0}\,\zeta_j\,\frac{1 + \alpha h_k^{\beta}}{v_k \, \sigma_{ok}} .
\label{eq:ptime}
\end{equation}
The quantity ratio \(q_j/q_0\) and the size factor \(\zeta_j \in \{0.88, 1.00, 1.16\}\) scale the batch. The machine speed \(v_k\) and the operator skill \(\sigma_{ok}\) scale the rate. The factor \(1 + \alpha h_k^{\beta}\) is the status effect, with \(\alpha = 0.85\) and \(\beta = 1.15\). The exponent above one makes the last part of the wear more expensive than the first, which is the empirical shape of a knife and of a needle that has started to heat. At \(h_k = 0\) the factor is one and the formula collapses to an ordinary flexible-job-shop time. At the health values used in the case, the factor is about \(1.45\) on the worn automatic cutter and about \(1.15\) on the healthy manual cutter. Those two numbers are computed in the worked example at the end of this section.

Equation~\eqref{eq:ptime} is the first place the twin enters the scheduler. A policy that treats \(p\) as a constant stored in a routing table is using \(\alpha = 0\), whether or not a screen on the wall shows \(h_k\).

\subsection{The heterogeneous graph of one decision}

At time \(t\) the ready set \(R(t)\) contains every operation whose predecessor in the same order is finished and whose order has been released. A candidate is a triple \((O_{ij}, k, o)\) with \(O_{ij} \in R(t)\), machine \(k\) eligible and free, operator \(o\) skilled and free, and the required token available. Candidates that fail any of those tests are not scored. Masking them is part of the model, not a cleanup after the network.

Each candidate induces a four-node heterogeneous graph. The operation node carries the time and due-date pressure of this bundle. The machine node carries health, load, energy, and failure. The operator node carries the skill gap. The shop node carries contention for the colour, the die, and the machine group. Edges are typed. A precedence edge is not drawn inside the candidate graph, because precedence has already been enforced by membership of \(R(t)\). The typed edges that remain are eligibility (operation--machine), assignment (operation--operator), seizure (operation--token, summarised on the shop node), and context (each node--shop). Figure~\ref{fig:rch-graph} draws one candidate.

\begin{figure}[ht]
\centering
\begin{tikzpicture}[
  n/.style={draw, circle, minimum size=1.15cm, font=\small, align=center},
  arr/.style={-{Stealth}, thick}
]
\node[n] (op) {op};
\node[n, right=2.2cm of op] (m) {machine};
\node[n, below=1.5cm of op] (o) {operator};
\node[n, below=1.5cm of m] (s) {shop};
\draw[arr] (op) -- node[above, font=\scriptsize] {eligible} (m);
\draw[arr] (op) -- node[left, font=\scriptsize] {skill} (o);
\draw[arr] (m) -- node[right, font=\scriptsize] {load} (s);
\draw[arr] (o) -- node[below, font=\scriptsize] {pool} (s);
\draw[arr] (op) -- node[above, font=\scriptsize, sloped] {token} (s);
\end{tikzpicture}
\caption{One-hop heterogeneous graph of a single candidate bundle. The full shop graph is the union of these candidate graphs at the current ready set. Attention pools the four nodes; it does not mix two candidates before the scores are compared.}
\label{fig:rch-graph}
\end{figure}

The twelve raw costs of a candidate, all scaled to lie near the unit interval, are
\begin{align}
\rho_1 &= p/45, &
\rho_2 &= s/11, &
\rho_3 &= \max(0,-\ell)/180, \label{eq:raw1}\\
\rho_4 &= (p + W)/ \max(30, d_j - t), &
\rho_5 &= h_k, &
\rho_6 &= \max(0, B_k - \bar B)/80, \label{eq:raw2}\\
\rho_7 &= E_{\mathrm{cycle}}, &
\rho_8 &= \phi, &
\rho_9 &= n^{\mathrm{col}} / n_{c_j}, \label{eq:raw3}\\
\rho_{10} &= n^{\mathrm{die}} / n_{\mathrm{die}}, &
\rho_{11} &= n^{\mathrm{group}}/3, &
\rho_{12} &= \max(0, 1.05 - \sigma_{ok}) . \label{eq:raw4}
\end{align}
Slack \(\ell\) is the due date minus the earliest plausible completion of the rest of the route. \(W\) is the base work remaining after this operation. The three contention counts \(n^{\mathrm{col}}\), \(n^{\mathrm{die}}\), and \(n^{\mathrm{group}}\) are the numbers of ready operations that want the same colour lot, the die pool, or the same machine group. They are the graph readout of the shop node: a mean, over neighbours, of demand. A candidate that looks short on \(\rho_1\) and disastrous on \(\rho_9\) is a khaki body started while two other khaki bodies are already waiting for the single lot.

The four node keys collapse the twelve costs along the node types,
\begin{align}
\kappa_{\mathrm{op}} &= \rho_1 + \rho_2 + 1.4\rho_3 + \rho_4 , \label{eq:key-op}\\
\kappa_{\mathrm{m}} &= \rho_5 + \rho_6 + 0.35\rho_7 + 1.6\rho_8 , \label{eq:key-m}\\
\kappa_{\mathrm{o}} &= \rho_{12} + 0.3\rho_{11} , \label{eq:key-o}\\
\kappa_{\mathrm{s}} &= \rho_9 + \rho_{10} + 0.5\rho_{11} . \label{eq:key-s}
\end{align}
Positive learned weights \(w_1,\ldots,w_{12}\) form the node costs
\begin{align}
c_{\mathrm{op}} &= w_1\rho_1 + w_2\rho_2 + w_3\rho_3 + w_4\rho_4 , \label{eq:node-op}\\
c_{\mathrm{m}} &= w_5\rho_5 + w_6\rho_6 + w_7\rho_7 + w_8\rho_8 , \label{eq:node-m}\\
c_{\mathrm{o}} &= w_{12}\rho_{12} , \label{eq:node-o}\\
c_{\mathrm{s}} &= w_8'\rho_9 + w_9\rho_{10} + w_{10}\rho_{11} , \label{eq:node-s}
\end{align}
where the shop-node weights are \(w_8, w_9, w_{10}\) in the parameter vector (the prime in \eqref{eq:node-s} is only there to keep the failure weight \(w_8\) of \eqref{eq:node-m} visually distinct; the implementation stores twelve weights and does not share them). Positivity is obtained by storing the logarithm and exponentiating, so a coordinate descent cannot flip a cost into a reward by accident.

\subsection{Attention and the preference}

Let \(a_{\mathrm{op}}, a_{\mathrm{m}}, a_{\mathrm{o}}, a_{\mathrm{s}}\) be attention biases and let \(g \in \R^4\) be a gate vector. The preference opens each node by
\begin{align}
\pi_{\mathrm{op}} &= 0.55 + g_1(\omega_1 + \omega_2) , \label{eq:open-op}\\
\pi_{\mathrm{m}} &= 0.55 + g_2(\omega_3 + 0.5\omega_1) , \label{eq:open-m}\\
\pi_{\mathrm{o}} &= 0.45 + g_3 \omega_4 , \label{eq:open-o}\\
\pi_{\mathrm{s}} &= 0.45 + g_4(\omega_2 + 0.25) . \label{eq:open-s}
\end{align}
Time and due-date weight open the operation node. Energy, and part of the makespan weight, open the machine node, because health and failure destroy makespan as well as kilowatt-hours. Balance opens the operator node. Due-date weight also opens the shop node, because token contention is how a due date is missed when the machine itself is free. The attention weights are a softmax over the opened keys,
\begin{equation}
\alpha_u = \frac{\exp(a_u + \pi_u \kappa_u)}{\sum_{u' \in \{\mathrm{op,m,o,s}\}} \exp(a_{u'} + \pi_{u'} \kappa_{u'})} .
\label{eq:alpha}
\end{equation}
The candidate score, which the policy minimises, is the opened, attention-weighted node cost
\begin{equation}
\psi(a) = \sum_u \alpha_u \, \pi_u \, c_u .
\label{eq:psi}
\end{equation}
Equation~\eqref{eq:alpha} is one-head, one-hop graph attention \cite{velickovic2018}. It is the operational policy. The multi-layer form used as the scalable model replaces the hand-grouped keys by learned maps. For each relation \(r\) and layer \(\ell\),
\begin{equation}
e_{uv}^{(r,\ell)} = \mathrm{LeakyReLU}\!\left(b_r^{\top}\big[W_r^{(\ell)} z_u^{(\ell)} \,\|\, W_r^{(\ell)} z_v^{(\ell)}\big]\right) ,
\label{eq:gat-full}
\end{equation}
\begin{equation}
\alpha_{uv}^{(r,\ell)} = \frac{\exp(e_{uv}^{(r,\ell)})}{\sum_{v' \in N_r(u)} \exp(e_{uv'}^{(r,\ell)})} ,
\qquad
z_u^{(\ell+1)} = \sigma\!\left(\sum_r \sum_{v \in N_r(u)} \alpha_{uv}^{(r,\ell)} W_r^{(\ell)} z_v^{(\ell)}\right) .
\label{eq:gat-update}
\end{equation}
The candidate score of the full model is a linear map of the concatenated final embeddings of the four nodes, again gated by \(\omega\). Equation~\eqref{eq:psi} is what Chapter~\ref{ch:estimation} runs. Equations~\eqref{eq:gat-full} and \eqref{eq:gat-update} are what a larger plant would train with the proximal policy method below, once more than a few dozen order books exist as logged transitions. Both are the same technology: a preference-conditioned attention pool over a resource-typed graph, followed by a masked choice of bundle.

\subsection{Markov decision process}

The dispatch process is a Markov decision process \(( \mathcal{S}, \mathcal{A}, P, R, \gamma )\). A state \(s_t \in \mathcal{S}\) contains the ready set, the free times and previous colours of machines and operators, the residual token capacities, the health vector, and the preference \(\omega\). An action \(a_t\) is a feasible bundle, or a bundle preceded by preventive maintenance when Section~\ref{sec:ssdt} says so. The transition \(P\) advances the clock to the next moment at which a machine and an operator are free, releases tokens of finished work, inserts any order whose release time has arrived, and updates health by the wear map of Section~\ref{sec:ssdt}. The reward is the change in the scalar cost. With \(J_\omega\) evaluated on the partial schedule,
\begin{equation}
r_t = -\big( J_\omega(S_{t+1}) - J_\omega(S_t) \big) ,
\label{eq:reward}
\end{equation}
so that a step which lengthens the day, the tardiness, the energy, the imbalance, or the failure sum returns a negative reward. The discount \(\gamma \in (0,1)\) is kept close to one because a setup avoided now pays its tardiness much later. The estimation identifies coefficients by the total \(J_\omega\) of the finished day, which is the undiscounted, finite-horizon form of the same objective.

The policy \(\pi_\theta(a \mid s, \omega)\) is a softmax over the negative scores of the feasible candidates,
\begin{equation}
\pi_\theta(a \mid s, \omega) = \frac{\exp(-\psi_\theta(a)/\tau)}{\sum_{a' \in \mathcal{F}(s)} \exp(-\psi_\theta(a')/\tau)} ,
\label{eq:policy}
\end{equation}
where \(\mathcal{F}(s)\) is the masked feasible set and \(\tau\) is a temperature. At execution, \(\tau \to 0\) and the policy takes the minimiser of \(\psi\). During stochastic training, a positive temperature supplies the exploration that a cross-entropy search gets instead from its sampling distribution.

\subsection{Training the full model with proximal policy optimisation}

The full model is trained by proximal policy optimisation \cite{schulman2017}. Let \(\theta\) collect the maps \(W_r\), the attention vectors, the gates, and the value-function weights. Along a trajectory the probability ratio of a new policy to the policy that collected the data is
\begin{equation}
\eta_t(\theta) = \frac{\pi_\theta(a_t \mid s_t, \omega)}{\pi_{\theta_{\mathrm{old}}}(a_t \mid s_t, \omega)} .
\label{eq:ratio}
\end{equation}
The clipped surrogate is
\begin{equation}
L^{\mathrm{clip}}(\theta) = \E_t \Big[ \min\big( \eta_t(\theta) \hat A_t,\; \mathrm{clip}(\eta_t(\theta), 1-\epsilon, 1+\epsilon)\, \hat A_t \big) \Big] ,
\label{eq:ppo}
\end{equation}
with \(\epsilon\) typically \(0.2\). The advantage \(\hat A_t\) is the generalised advantage estimate of the scalar reward \eqref{eq:reward}. A value head \(V_\phi(s,\omega)\) is trained by squared error against the discounted return, and an entropy bonus \(\beta_H \mathcal{H}(\pi_\theta)\) keeps \eqref{eq:policy} from collapsing to a single bundle before the resource features have been used. The parameter step maximises \(L^{\mathrm{clip}} - c_V L^{\mathrm{value}} + \beta_H \mathcal{H}\).

Preference conditioning is trained by sampling \(\omega\) with each episode, not by training four unrelated networks and hoping their weights happen to be comparable. A single network that has seen many \(\omega\) is what Equation~\eqref{eq:open-op}--\eqref{eq:open-s} approximate when the gates \(g\) are identified at one operating preference and then reused. Chapter~\ref{ch:estimation} is explicit about which of these two regimes the numbers come from: the coefficients were identified at the balanced preference \(\omega = (0.34, 0.33, 0.18, 0.15)\), and the gate formulas were not re-fit for every corner of the simplex. The tolerance sweep of the twin, not the gate formulas, is the knob the estimation shows moving cleanly.

\subsection{Identification used in the estimation system}

The executable score has twenty coefficients: twelve log-weights, four attention biases, and four gates. Direct policy-gradient training of a deep network is not required to identify twenty numbers on a shop of this size, and it would make the tables depend on an optimiser trajectory that a reader cannot see. The estimation therefore uses the cross-entropy method. A Gaussian \(\mathcal{N}(\mu, \mathrm{diag}(\sigma^2))\) proposes a population of coefficient vectors. Each vector is scored by the mean of \(J_\omega\) on a fixed set of training order books, always relative to first-in-first-out on the same book. The elite fraction updates \(\mu\) and \(\sigma\). The search starts from a logarithmic prior that already resembles a blend of shortest processing time, slack, and health, so the first generation is a scheduler rather than a random permutation. Algorithm~\ref{alg:cem} states the loop. The training curve is reported in Chapter~\ref{ch:estimation}.

\begin{algorithm}[ht]
\caption{Cross-entropy identification of the RCH-GAT score}
\label{alg:cem}
\begin{algorithmic}[1]
\State Initialise \(\mu\) from the short-time, slack, and health prior; set \(\sigma \leftarrow 0.35\)
\For{generation \(= 1, \ldots, G\)}
  \For{sample \(= 1, \ldots, N\)}
    \State Draw \(\theta \sim \mathcal{N}(\mu, \mathrm{diag}(\sigma^2))\)
    \State For each training book, run Algorithm~\ref{alg:dispatch} and record \(J_\omega\)
    \State Set the fitness to the mean \(J_\omega\) across the training books
  \EndFor
  \State Keep the elite fraction of \(\theta\)
  \State Set \(\mu\) to their mean and \(\sigma\) to their standard deviation, floored by a small positive constant
\EndFor
\State \Return the elite \(\theta\) with the best fitness observed
\end{algorithmic}
\end{algorithm}

\begin{algorithm}[ht]
\caption{Online dispatch with a masked resource bundle}
\label{alg:dispatch}
\begin{algorithmic}[1]
\State \(t \leftarrow 0\); mark every operation unfinished; set health to the measured \(h_k(0)\)
\While{some operation is unfinished}
  \State Release any order with \(r_j \le t\), and release tokens of occupations that have ended
  \State Build the feasible set \(\mathcal{F}\) of bundles at time \(t\)
  \If{\(\mathcal{F}\) is not empty}
    \State Score each bundle by \eqref{eq:psi}, applying the maintenance test of Algorithm~\ref{alg:ssdt}
    \State Start the minimiser; update free times, colour memory, tokens, health, energy, and \(\phi\)
  \Else
    \State Advance \(t\) to the next completion or the next release
  \EndIf
\EndWhile
\end{algorithmic}
\end{algorithm}

\subsection{What the mask guarantees}

\begin{proposition}
\label{prop:mask}
Suppose every order has a finite route, every operation has at least one eligible machine group, every machine group has at least one operator with positive skill, and every token pool that an operation requires has capacity at least one. Then Algorithm~\ref{alg:dispatch} finishes every released order in finite time, and every started bundle satisfies precedence, machine capacity, operator capacity, and token capacity.
\end{proposition}

\begin{proof}
Every started bundle is taken from \(\mathcal{F}\), and \(\mathcal{F}\) contains only bundles whose predecessor is finished, whose machine and operator are free, whose skill is positive, and whose token pool has a free unit. Capacity and precedence therefore hold by construction of the mask. A started bundle occupies its machine, operator, and token for a finite duration, because base times, size factors, and health factors are finite and health is clipped. At the occupation's end the three resources are released. If \(\mathcal{F}\) is empty and work remains, either a resource is occupied, in which case the clock advances to a release, or an order is unreleased, in which case the clock advances to its release time. The state cannot cycle: each advance of the clock to a completion finishes one operation, and there are finitely many operations. Hence the loop ends with every operation finished.
\end{proof}

The proposition is about feasibility, not about optimality. A first-in-first-out policy satisfies it and still produces the poor tardiness reported in Chapter~\ref{ch:estimation}. The proposition is what makes the comparison in that chapter a comparison of decisions rather than a comparison of repair heuristics. Every method in the estimation, including the classical rules, dispatches through the same mask. They differ only in the score that ranks \(\mathcal{F}\).

\subsection{A numerical assignment without the twin's veto}

Consider one ready operation, the shell cut of a navy backpack, quantity \(80\), size M, so that \(q_j/q_0 = 1\) and \(\zeta_j = 1\). The base time is \(\tau = 24\) minutes. Cutter-A has speed \(v = 1.28\), power \(6.8\,\mathrm{kW}\), wear rate \(0.00115\) per minute, and health \(h = 0.58\). Cutter-B has speed \(0.74\), power \(2.4\,\mathrm{kW}\), and health \(h = 0.22\). Operator Cut-1, with skill \(1\), is free for either machine. No setup is charged in this comparison. Equation~\eqref{eq:ptime} on Cutter-A gives
\begin{equation}
1 + \alpha h^{\beta} = 1 + 0.85 \times 0.58^{1.15} = 1.4543 ,
\end{equation}
\begin{equation}
p_A = 24 \times 1.4543 / 1.28 = 27.268 \ \mathrm{minutes} .
\end{equation}
On Cutter-B,
\begin{equation}
1 + \alpha h^{\beta} = 1 + 0.85 \times 0.22^{1.15} = 1.1490 ,
\end{equation}
\begin{equation}
p_B = 24 \times 1.1490 / 0.74 = 37.265 \ \mathrm{minutes} .
\end{equation}
A shortest-processing-time rule sees \(27.268 < 37.265\) and assigns Cutter-A. The energy of that cycle, from the first term of \eqref{eq:f3} and before the failure term, is
\begin{equation}
E_A = 6.8 \times \frac{27.268}{60} \times (1 + 0.45 \times 0.58) = 3.896\ \mathrm{kWh} ,
\end{equation}
against
\begin{equation}
E_B = 2.4 \times \frac{37.265}{60} \times (1 + 0.45 \times 0.22) = 1.638\ \mathrm{kWh} .
\end{equation}
The fast machine is more than twice as hungry. The failure probability, computed in Section~\ref{sec:ssdt} from the hazard at these health values, is \(0.106\) on A and \(0\) on B. Expected repair delay on A is \(0.106 \times 48 = 5.09\) minutes, so the occupation that shortest processing time thought was \(27.3\) minutes is \(32.4\) minutes once the risk is priced, and it leaves Cutter-A at a predicted health of \(0.58 + 0.00115 \times 27.268 = 0.611\).

The attention score does not have to choose B every time. If the due date is already broken and the other cutter is the bottleneck of a later order, paying \(3.9\,\mathrm{kWh}\) on A can be the lower score under a time-heavy \(\omega\). The point of the numbers is narrower. Any rule that ranks on nominal or even on health-scaled processing time alone will take A. Any score that includes \(\rho_5\), \(\rho_7\), and \(\rho_8\) is allowed to refuse A. That refusal is the resource-and-status content of \(\psi\). Whether the refusal should instead be a sixteen-minute maintenance of A is the twin's decision, made with the same candidate in the next section.

\subsection{Attention arithmetic on the two cutters}

The preceding subsection stopped at the physical times. This one applies the identified score, so that Equation~\eqref{eq:psi} is a number rather than a shape. The coefficient vector is the one stored by Algorithm~\ref{alg:cem} and used for every table in Chapter~\ref{ch:estimation}. Exponentiating the twelve log-weights gives
\begin{equation}
\begin{aligned}
w = (\, &0.288,\ 1.029,\ 0.731,\ 0.235,\ 0.606,\ 5.660,\\
&0.912,\ 1.269,\ 0.133,\ 0.917,\ 0.353,\ 1.657 \,).
\end{aligned}
\label{eq:wstar}
\end{equation}
The order is processing time, setup, slack penalty, critical ratio, health, load deviation, energy, failure probability, colour contention, die contention, group contention, and skill gap. Two facts are already visible. The load weight \(5.660\) is several times any other weight, so a machine that has already run far ahead of its neighbours is expensive even when it is healthy. The colour-contention weight \(0.133\) is the smallest of the twelve, which is the coefficient-level reason the resource ablation in Chapter~\ref{ch:estimation} moves the day only a little.

At the balanced preference \(\omega = (0.34, 0.33, 0.18, 0.15)\) the openings \eqref{eq:open-op}--\eqref{eq:open-s} evaluate to
\begin{equation}
\pi = (1.121,\ 0.908,\ 0.734,\ -0.490).
\label{eq:pi-num}
\end{equation}
The shop-node opening is negative. The identified gate on that node learned to shrink, rather than enlarge, the contention term under this preference. A negative opening does not delete the node. Multiplied into the key, it reduces attention when contention is high; multiplied into the cost, it can even change the sign of that node's contribution. In the cutter arithmetic below the attention mass on the shop node is about \(0.03\), so the sign changes the score by roughly a hundredth. Token pressure is a weak voice, not an inverted one in any way that moves a schedule. That is an empirical outcome of the identification, reported here so that the ablation is not a surprise.

Construct two candidates for the navy shell, both with Cut-1 free, no setup, no slack penalty, and the same token pressure: one other cutting operation is also ready, and the navy pool is half used. Critical ratios are \(0.80\) on A and \(0.95\) on B, because B's longer cycle eats more of the remaining due-date margin. Load deviation is \(0\) on A and \(0.10\) on B. The twelve raw costs are then
\begin{align}
\rho^{A} &= (0.606,\ 0,\ 0,\ 0.800,\ 0.580,\ 0,\ 3.896,\ 0.106,\ 0.500,\ 0.500,\ 0.667,\ 0), \\
\rho^{B} &= (0.828,\ 0,\ 0,\ 0.950,\ 0.220,\ 0.100,\ 1.638,\ 0,\ 0.500,\ 0.500,\ 0.667,\ 0).
\end{align}
Node keys and node costs follow \eqref{eq:key-op}--\eqref{eq:node-s}. For Cutter-A,
\begin{equation}
\kappa^{A} = (1.406,\ 2.113,\ 0.200,\ 1.333), \qquad
c^{A} = (0.363,\ 4.038,\ 0,\ 0.760).
\end{equation}
The machine-node cost \(4.038\) is almost entirely energy: \(0.912 \times 3.896 = 3.55\), with health and failure making up the rest. For Cutter-B,
\begin{equation}
\kappa^{B} = (1.778,\ 0.893,\ 0.200,\ 1.333), \qquad
c^{B} = (0.462,\ 2.193,\ 0,\ 0.760).
\end{equation}
Softmax attention \eqref{eq:alpha} with the biases \(a = (1.082, 0.969, -0.707, 0.672)\) puts mass
\begin{equation}
\alpha^{A} = (0.422,\ 0.531,\ 0.017,\ 0.030), \qquad
\alpha^{B} = (0.742,\ 0.203,\ 0.020,\ 0.035).
\end{equation}
On the worn cutter, attention sits on the machine node, because that node's key is large. On the healthy cutter, attention sits on the operation node, because the machine key has collapsed. The scores are
\begin{equation}
\psi(A) = 2.108, \qquad \psi(B) = 0.776.
\end{equation}
The policy takes B. Shortest processing time, using \(27.3 < 37.3\), takes A. The difference is not a hidden rule. It is the machine-node cost of \(6.8\,\mathrm{kW}\) at health \(0.58\), amplified by an attention weight of \(0.531\).

Table~\ref{tab:weights} is the same vector, written so a reader can see which physical effect survived identification.

\begin{table}[ht]
\centering
\caption{Identified positive weights on the twelve costs, from \eqref{eq:wstar}. A large weight means the identification was willing to avoid that cost even when other costs rose.}
\label{tab:weights}
\small
\begin{tabular}{lrll}
\toprule
Cost & Weight & Node & What a large value punishes \\
\midrule
Processing time & 0.288 & operation & a long cycle \\
Setup & 1.029 & operation & a colour change \\
Slack penalty & 0.731 & operation & an already broken due date \\
Critical ratio & 0.235 & operation & a route that still has a lot left \\
Health & 0.606 & machine & a worn tool \\
Load deviation & 5.660 & machine & a machine far ahead of the others \\
Energy & 0.912 & machine & a hungry cycle \\
Failure probability & 1.269 & machine & a likely interruption \\
Colour contention & 0.133 & shop & a crowded thread lot \\
Die contention & 0.917 & shop & a crowded die pool \\
Group contention & 0.353 & shop & a crowded machine group \\
Skill gap & 1.657 & operator & a weak pairing \\
\bottomrule
\end{tabular}
\end{table}

The ranking inside the table is the design story of the first technology after it has been fit, rather than before. Balance of machine load dominates, which is why imbalance falls in every scenario of Chapter~\ref{ch:estimation}. Skill gap is the next strongest single weight, so the operator node is not decorative: a slow pairing is refused when a qualified operator is free. Failure probability outweighs raw processing time. Colour contention, the feature that sounds most like a school-bag factory, is almost silent. Die contention is not silent. The two token types are not interchangeable, and a model that had used one ``resource'' weight for both would have blurred a real difference. Setups are punished about as strongly as energy. A navy-to-red change of eleven minutes is allowed to break a tie between two otherwise similar sewing machines.

\subsection{A strap, two machines, and one missing skill}

The cutter example has a free specialist. Sewing often does not. Consider the strap of the same backpack. The base time is \(30\) minutes. Bartack-1 has speed \(1.00\), power \(1.15\,\mathrm{kW}\), and health \(0.34\). Lockstitch-1 has speed \(1.05\), power \(0.75\,\mathrm{kW}\), and health \(0.24\). Sew-2 is the bartack specialist, skill \(1.08\) on the heavy machine and \(0.96\) on a lockstitch. Sew-3 can run only a lockstitch, at skill \(1.00\). Sew-1 can run the bartack at skill \(0.78\), which is legal and poor.

Equation~\eqref{eq:ptime} at these health values, for Sew-2 on the bartack, has health factor \(1 + 0.85 \times 0.34^{1.15}\). The factor is about \(1.25\), and the processing time is about \(30 \times 1.25 / 1.08 \approx 34.7\) minutes. The same operator on Lockstitch-1 is slower in skill and faster in machine speed, and the health factor is closer to \(1.16\), giving a cycle near \(30 \times 1.16 / (1.05 \times 0.96) \approx 34.5\) minutes. The two cycles are almost the same length. Shortest processing time will treat them as a tie and then break the tie by setup or by order identity. The score will not, because the skill gap, the load, and the failure probability are different features.

Sew-3 on the bartack is not a candidate. The mask of Proposition~\ref{prop:mask} removes it before \(\psi\) is computed. This is the dual-resource constraint in one line \cite{xu2011}. A chromosome that picks ``bartack'' and then looks for an operator afterwards can discover the infeasibility only in the decoder. The graph never scores the infeasible pair, so the attention cannot waste mass on it.

If Sew-2 is already inside a trolley shell, the only legal bartack operator left is Sew-1 at skill \(0.78\). The skill-gap cost is \(\max(0, 1.05 - 0.78) = 0.27\), and the weight on that cost is \(1.657\), so the operator node alone contributes about \(0.45\) before attention and opening. The cycle also lengthens by the ratio \(1.08/0.78\). The policy will often leave the strap on a lockstitch rather than staff the bartack with the weak pairing, unless the lockstitches are the congested group. That choice is exactly what the operator node was added to make. A machine-only graph of the kind in Song, Chen, Li, and Cao \cite{song2023} would assign the bartack whenever it is free, and the simulator would then have to repair the assignment by waiting for Sew-2. The wait is a hidden setup. Here it is a term in \(\psi\).

\subsection{How a preference would have to be re-fit}

Equation~\eqref{eq:pi-num} is not a family of policies. It is one policy, at one \(\omega\), whose gates were identified jointly with the weights. Changing \(\omega\) in the opening formulas without repeating Algorithm~\ref{alg:cem} does not trace a reliable front, and Chapter~\ref{ch:estimation} does not claim that it does. The correct way to ask for an energy-heavy shift is to change the scalar \eqref{eq:scalar} and identify again. On this cell that identification is a few dozen generations of a twenty-dimensional Gaussian, which is cheap beside a day of production and is still a different vector from \eqref{eq:wstar}.

What the openings do provide, even at a single fit, is an audit. A reader can see that the shop node was driven toward silence, that the machine node stayed open, and that the operation node carries makespan and due date together. If a later fit at a due-date-heavy \(\omega\) raised the colour weight from \(0.133\) and flipped the shop gate from negative to positive, the resource ablation would be expected to grow. Until that fit is run, the honest description of RCH-GAT on this cell is: load, skill, failure, setup, and energy are the active costs; colour contention is specified and nearly inactive.

\subsection{Scaling, and what would break it}

The divisors in \eqref{eq:raw1}--\eqref{eq:raw4} are not natural constants. Processing time is divided by \(45\) minutes because a typical sewing or cutting cycle on this cell lies between \(15\) and \(60\). Slack is divided by \(180\) minutes because a due-date miss of half a shift should be the same order of magnitude as a processing-time cost near one. Energy is left in kilowatt-hours because a cutter cycle is already a few units and a sewing cycle is a fraction of a unit. If the cell were rescaled to single bags instead of batches of eighty, every divisor would be wrong and the weights in \eqref{eq:wstar} would have to be identified again. The graph and the mask would not. This is the usual separation between a representation and a calibration. The thesis treats the divisors as part of the calibration of this cell, and it does not present them as universal.

A second breakage is a health indicator that is not on the unit interval. Equation~\eqref{eq:ptime} with \(\alpha = 0.85\) assumes \(h\) near \(0\) is a new tool and \(h\) near \(1\) is a failed one. A raw motor current must be mapped onto that interval before it enters \(\rho_5\). The map is the feature extraction reviewed by Lei and co-authors \cite{lei2018}. Putting the current in directly would make the health weight meaningless and would make the chance constraint fire on the wrong cycles.

A third breakage is a token pool whose capacity changes during the day, for example a khaki lot that is discovered to be short at noon. The shop node can represent the new capacity immediately, because \(n_c\) is read at the decision epoch and is not compiled into the weights. The weights do not have to be identified again for a shortage. The shortage scenario of Chapter~\ref{ch:estimation} uses exactly that property. The small movement in that scenario is then not a failure of the data path. It is a consequence of the small colour weight.

\section{Status-synchronized chance constraints}
\label{sec:ssdt}

\subsection{What is synchronised}

The twin of machine \(k\) is the scalar \(h_k(t)\) together with the two maps that the dispatcher reads: the duration factor in Equation~\eqref{eq:ptime} and the failure probability \(\phi\) below. Synchronisation means that \(h_k\) used at a dispatch epoch is the health implied by the work the machine has already done and by the latest accepted sensor reading, not the health written on the morning sheet. In the estimation plant the simulator is the sensor, so the wear recursion is exact. In a physical cell the recursion is the drift, and a Kalman update absorbs the residual of the sensor. Both are specified, and the chapter does not pretend that the estimation ran a noisy sensor.

\subsection{Drift, diffusion, and the sensor}

Between decisions the design model is the Wiener process reviewed in Section~\ref{sec:tech2-review},
\begin{equation}
\mathrm{d}h_k = \eta_k \, \mathrm{d}t + \sigma_k \, \mathrm{d}B_t ,
\label{eq:sde}
\end{equation}
with a machine-specific wear rate \(\eta_k\). Over a cycle of processing length \(p\), the increment has mean \(\eta_k p\) and variance \(\sigma_k^2 p\). A sensor does not observe \(h_k\). It observes
\begin{equation}
y_k = h_k + \epsilon_k , \qquad \epsilon_k \sim \mathcal{N}(0, \rho_k^2) ,
\label{eq:sensor}
\end{equation}
where \(\rho_k\) is the noise of the health indicator after the feature extraction that Lei and co-authors survey \cite{lei2018}. The indicator itself, whether it is a current, a temperature, or a thread-break rate, is assumed already reduced to the unit interval on which \(h_k = 0\) is a new tool and \(h_k = 1\) is a tool the maintenance standard calls failed.

The Kalman filter is the revision of this linear-Gaussian pair \cite{kalman1960}. Predict, over a cycle just completed,
\begin{equation}
\hat h^- = \hat h + \eta_k p , \qquad v^- = v + \sigma_k^2 p .
\label{eq:kalman-pred}
\end{equation}
Update, when the reading \(y\) arrives,
\begin{equation}
K = \frac{v^-}{v^- + \rho_k^2} , \qquad \hat h = \hat h^- + K(y - \hat h^-) , \qquad v = (1-K)v^- .
\label{eq:kalman-upd}
\end{equation}
The dispatcher uses \(\hat h\), not \(y\). A single noisy spike therefore cannot throw the cutter into maintenance, and a slow drift cannot hide behind one optimistic reading. Particle filters replace \eqref{eq:kalman-pred}--\eqref{eq:kalman-upd} if a later plant shows that the residual is heavy-tailed \cite{arulampalam2002}. The estimation does not need that replacement, because its plant applies the drift directly:
\begin{equation}
h_k \leftarrow \min\big(1.25,\; h_k + \eta_k p \big) ,
\label{eq:plant-wear}
\end{equation}
with \(\eta_k\) equal to the wear coefficient of the machine (for the automatic cutter, \(0.00115\) per minute). Equation~\eqref{eq:plant-wear} is \eqref{eq:kalman-pred} with the posterior variance held at zero. Reporting it this way keeps the tables reproducible and keeps the Kalman equations available for the cell in which the variance is not zero.

\subsection{Failure probability as a hazard}

Given a health at the start of a cycle, the intensity used by the dispatcher is
\begin{equation}
\lambda(h, p) = \lambda_0 \big(\max(0, h_{\mathrm{mid}} - h_{\mathrm{fail}})\big)^{\gamma} ,
\label{eq:hazard}
\end{equation}
with \(\lambda_0 = 2.4\), \(h_{\mathrm{fail}} = 0.40\), \(\gamma = 1.35\), and the mid-cycle health
\begin{equation}
h_{\mathrm{mid}} = h + \tfrac{1}{2} \eta_0 p .
\label{eq:mid}
\end{equation}
The constant \(\eta_0 = 0.0004\) in the mid-cycle shift is a small look-ahead, not the full machine wear rate. The failure probability of a cycle of \(p\) minutes is the exponential formula for a constant intensity,
\begin{equation}
\phi(h, p) = 1 - \exp\!\big(-\lambda(h,p) \, p / 60\big) .
\label{eq:phi}
\end{equation}
Below the threshold \(h_{\mathrm{fail}}\) the intensity is zero and \(\phi = 0\). Above it, intensity grows faster than linearly because \(\gamma > 1\). This is the operational chance model. It is monotone in health and in cycle length, it is cheap to evaluate inside the inner loop of Algorithm~\ref{alg:dispatch}, and it is the \(\phi\) that Chapter~\ref{ch:estimation} sums into \(f_5\).

The Gaussian derivation that justifies a threshold of this shape starts from \eqref{eq:sde}. The probability of crossing a ceiling \(h^{\star}\) during a cycle is
\begin{equation}
\Prob(h + \eta p + \sigma \sqrt{p}\, Z \ge h^{\star}) = 1 - \Phi\!\left( \frac{h^{\star} - h - \eta p}{\sigma \sqrt{p}} \right) ,
\label{eq:gauss-cross}
\end{equation}
where \(Z\) is standard normal and \(\Phi\) is its distribution function. Requiring \eqref{eq:gauss-cross} to be at most \(\varepsilon\) is equivalent to
\begin{equation}
h + \eta p + \sigma \sqrt{p}\, \Phi^{-1}(1-\varepsilon) \le h^{\star} .
\label{eq:gauss-threshold}
\end{equation}
Equation~\eqref{eq:gauss-threshold} is the classical chance constraint \cite{charnes1959} specialised to a Wiener increment. The hazard \eqref{eq:phi} is the form used when the plant's interruption process is defined by an intensity rather than by a single Gaussian crossing, which is the case in the estimation. Both forms give the dispatcher the same kind of object: a monotone function of \((h,p,\varepsilon)\) that can veto an assignment. The estimation's veto is \eqref{eq:phi}, compared with \(\varepsilon\), together with a hard ceiling on the predicted health at the end of the cycle,
\begin{equation}
h + \eta_k p \le h^{\mathrm{ceil}} = 0.90 .
\label{eq:ceil}
\end{equation}
If either \eqref{eq:phi} exceeds \(\varepsilon\) or \eqref{eq:ceil} fails, the bundle is not started as it stands.

\subsection{The maintenance action}

A vetoed bundle is not deleted. Deleting it would remove the fast cutter from \(\mathcal{F}\) exactly when the order book most needs a cutter, and Proposition~\ref{prop:mask} would then depend on the second cutter happening to be free. The twin instead offers the same bundle with a preventive maintenance prefixed,
\begin{equation}
h_k \leftarrow 0.30\, h_k + 0.04 , \qquad \Delta \leftarrow \Delta + \Delta^{\mathrm{pm}} ,
\label{eq:pm-map}
\end{equation}
with \(\Delta^{\mathrm{pm}} = 16\) minutes. The processing time and the failure probability are recomputed at the renewed health. The maintenance occupies the machine and the operator. It does not seize the die or the thread lot: the tool change on a cutter is not a fabric-die change. Energy of the combined occupation uses the renewed health, because the renewed machine is the one that runs.

The coefficient \(0.30\) is a partial renewal, not a new machine. A knife change or a needle change and a short check bring the indicator most of the way down and leave a small residual, the \(0.04\). This is the opposite of an as-good-as-new assumption. It is also the reason maintenance remains attractive on a worn cutter and useless on a healthy one: if \(h\) is already \(0.10\), Equation~\eqref{eq:pm-map} barely moves it, \(\phi\) was already zero, and the veto does not fire.

The score of a maintained bundle is the score of the renewed cycle plus a maintenance penalty proportional to the machine-node opening \(\pi_{\mathrm{m}}\), minus a credit for the failure and health costs that the renewal removes. The penalty keeps the policy from maintaining a machine whose \(\phi\) exceeds \(\varepsilon\) by a negligible amount when another feasible bundle is clean. The credit keeps the policy from refusing a renewal that removes a large \(\rho_8\). Both terms are part of \(\psi\), so maintenance is not an external rule glued on after the attention. It is a second version of the same candidate.

\subsection{When the decision is taken again}

Because Algorithm~\ref{alg:dispatch} rebuilds \(\mathcal{F}\) whenever a resource is free, the twin does not need a separate rescheduling solver. A health value that changed at the last completion is already inside \(\rho_5\) and \(\rho_8\) of the next score. A rush order, once its release time is reached, enters \(R(t)\) and therefore enters \(\mathcal{F}\). This is event-triggered rescheduling in the operational sense: the trigger is ``a resource became free, or an order became released, or a maintenance ended,'' and the response is a new masked argmin of \(\psi\).

A heavier trigger is available and is not required for correctness. If a sensor residual \(|y - \hat h|\) from Equation~\eqref{eq:kalman-upd} exceeds a band, the cell may discard the ranking it has been following and recompute immediately, even in the middle of a planned idle. The estimation plant has no such residual, so its trigger is the resource-free event alone. The effect that the literature calls closed-loop rescheduling is still present. No priority list is frozen at \(t = 0\). The classical rules in the comparison are given the same courtesy: earliest due date, at time \(t\), sees the orders that have been released by \(t\), including a rush order after its release. The comparison does not handicap the baselines with a stale list.

\begin{algorithm}[ht]
\caption{Chance-constrained maintenance test inside a candidate}
\label{alg:ssdt}
\begin{algorithmic}[1]
\State Compute \(p\) by \eqref{eq:ptime} at the current \(h_k\), then \(\phi\) by \eqref{eq:phi}
\State Predict the ending health by \(h_k + \eta_k p\)
\If{\(\phi > \varepsilon\) or \(h_k + \eta_k p > 0.90\)}
  \State Apply \eqref{eq:pm-map}; recompute \(p\) and \(\phi\); charge \(\Delta^{\mathrm{pm}}\)
  \State Mark the candidate as maintained
\EndIf
\State Return the occupation length, the energy, \(\phi\), and the maintained flag to the scorer \(\psi\)
\end{algorithmic}
\end{algorithm}

\subsection{Coupling the twin back into the graph}

The two technologies share nodes. The machine node of Figure~\ref{fig:rch-graph} is where the twin writes. After every completion, Equation~\eqref{eq:plant-wear} or the Kalman pair \eqref{eq:kalman-pred}--\eqref{eq:kalman-upd} replaces \(h_k\). The next time that machine appears in a candidate, \(\rho_5\), \(\rho_7\), and \(\rho_8\) are different numbers, so \(\kappa_{\mathrm{m}}\) and \(c_{\mathrm{m}}\) are different numbers, so the softmax \eqref{eq:alpha} can put more or less mass on the machine node, so \(\psi\) can change the chosen bundle. After a maintenance, the same write-back happens immediately, before the processing part of the occupation, which is why the processing time of a maintained cycle is the time at the renewed health.

There is a second, quieter coupling. Expected failure delay is added to the machine's free time. A policy that sends work to a worn machine does not only risk a future repair in the abstract. It blocks that machine, and the operator standing at it, for the expected repair. Subsequent candidates see a later free time. Status therefore changes contention, not only the row of the table that stores \(h_k\). This is why an ablation that zeroes the health features but still runs the plant's breakdown delay is not the same thing as an ablation that removes maintenance. Chapter~\ref{ch:estimation} runs both: one policy sees health and may maintain, one policy sees health and may not maintain, and one policy does not see health and may not maintain.

\begin{figure}[ht]
\centering
\begin{tikzpicture}[
  box/.style={draw, rounded corners, align=center, font=\small, minimum width=2.5cm, minimum height=0.85cm},
  arr/.style={-{Stealth}, thick}
]
\node[box] (h) {Health \(h_k\)};
\node[box, right=1.15cm of h] (phi) {\(\phi\) and \(p(h)\)};
\node[box, right=1.15cm of phi] (psi) {Score \(\psi\)};
\node[box, below=1.0cm of psi] (act) {Bundle or\\maintenance};
\node[box, below=1.0cm of h] (up) {Wear or\\Kalman update};
\draw[arr] (h) -- (phi);
\draw[arr] (phi) -- (psi);
\draw[arr] (psi) -- (act);
\draw[arr] (act) -- (up);
\draw[arr] (up) -- (h);
\end{tikzpicture}
\caption{Coupling of SSDT-CC into RCH-GAT. The score never reads a raw sensor. It reads a health state that maintenance and ordinary wear both rewrite.}
\label{fig:coupling}
\end{figure}

\subsection{The same cut, with the veto turned on}

Return to the navy shell of Section~\ref{sec:rch}. On Cutter-A, \(h = 0.58\) and \(p = 27.268\). The mid-cycle health in \eqref{eq:mid} is \(0.5855\). The intensity is
\begin{equation}
\lambda = 2.4 \times (0.5855 - 0.40)^{1.35} = 0.2468 \ \mathrm{per\ hour} ,
\end{equation}
and
\begin{equation}
\phi_A = 1 - \exp(-0.2468 \times 27.268/60) = 0.1061 .
\end{equation}
On Cutter-B, mid-cycle health \(0.227\) lies below \(0.40\), so \(\lambda = 0\) and \(\phi_B = 0\). The predicted ending health on A is \(0.611\), which does not exceed \(0.90\). The failure test does exceed a tolerance of \(\varepsilon = 0.10\). Algorithm~\ref{alg:ssdt} therefore maintains Cutter-A before the cut.

The renewal gives \(h' = 0.30 \times 0.58 + 0.04 = 0.214\). The health factor at \(h'\) is \(1.1443\), and the processing time falls to
\begin{equation}
p_A' = 24 \times 1.1443 / 1.28 = 21.456 \ \mathrm{minutes} .
\end{equation}
With the sixteen-minute maintenance the occupation is \(37.456\) minutes, and \(\phi\) at the renewed health is zero because \(0.214 < 0.40\). Energy of the maintained occupation is \(4.654\,\mathrm{kWh}\). Compare the three options on this single cycle, ignoring setup.

\begin{table}[ht]
\centering
\caption{One navy shell cut, three ways of treating Cutter-A. Times are minutes. Energy is kilowatt-hours. Failure delay is the expectation \(48\phi\).}
\label{tab:worked}
\begin{tabular}{lrrrr}
\toprule
Option & Process & Extra & Total & Energy \\
\midrule
Cutter-A, no maintenance & 27.268 & 5.09 & 32.36 & 3.90 + repair share \\
Cutter-A, maintenance first & 21.456 & 16.00 & 37.46 & 4.65 \\
Cutter-B, no maintenance & 37.265 & 0.00 & 37.27 & 1.64 \\
\bottomrule
\end{tabular}
\end{table}

Table~\ref{tab:worked} is the whole argument of the twin in one row. Maintenance is not a free improvement of the immediate cycle. It is slower than running the worn cutter and accepting the expected repair, and it uses more energy than the manual cutter. What it buys is a failure probability of zero on this cycle and a cutter whose next cycles run at health \(0.214\) rather than at a health that has stepped to \(0.611\). The day-level effect of that purchase cannot be seen in one row. It is the content of the wear scenario in Chapter~\ref{ch:estimation}, where the same sixteen-minute action, repeated when the test fires, cuts the expected failure count from about thirty to under one and pulls the makespan down with it.

A tolerance of \(0.10\) is a management choice. At \(\varepsilon = 0.20\) this particular cycle would not be vetoed by \(\phi\), because \(0.106 < 0.20\), and the ceiling \(0.90\) would not veto it either. The cutter would run worn. Section~\ref{sec:estimate} sweeps \(\varepsilon\) for that reason. The formula that decides is \eqref{eq:phi}, not a hidden list of machines that must be maintained.

\subsection{Feasibility of the maintained bundle}

Maintenance does not break Proposition~\ref{prop:mask}. The maintained bundle uses the same machine and the same operator, for a longer finite duration, and then releases them. Token capacity is unchanged because the token is not held during the tool change. The renewed processing time is finite. A shop in which every cutter is worn and every operator is free will maintain and then cut, not deadlock. A shop in which the only cutter operator is busy will not start the maintenance until that operator is free, because the candidate was never in \(\mathcal{F}\). The twin does not invent capacity.

There is a real cost that the proposition does not limit. If \(\varepsilon\) is set absurdly low, the test fires on mild health and the day fills with sixteen-minute renewals. The scalar \(J_\omega\) punishes the makespan of those renewals, and the estimation's sweep shows the resulting exchange between \(\phi\) and maintenance count. The technology includes \(\varepsilon\) as a visible parameter because hiding it inside the network would make that exchange impossible to audit.

\subsection{Three cycles of the same cutter}

One veto does not explain a day. Three identical navy shell cuts, back to back on Cutter-A, with Cut-1 free and no setup, show the horizon the twin is actually buying. Without maintenance the plant follows Equation~\eqref{eq:plant-wear}. With the tolerance at \(0.10\), the first cycle is vetoed, the tool is renewed once, and the next three cycles run at the renewed drift. Table~\ref{tab:threecycle} uses the same formulas as the single-cycle example.

\begin{table}[ht]
\centering
\caption{Three successive shell cuts on Cutter-A. ``Charged'' time is processing plus expected repair delay, or processing plus the single sixteen-minute maintenance on the renewed path. Health is the value at the start of the cycle.}
\label{tab:threecycle}
\small
\begin{tabular}{lrrrr}
\toprule
Path & Cycle & Start \(h\) & Process (min) & Charged (min) \\
\midrule
No maintenance & 1 & 0.580 & 27.27 & 32.36 \\
No maintenance & 2 & 0.611 & 27.80 & 34.12 \\
No maintenance & 3 & 0.643 & 28.35 & 35.99 \\
Renewed once & 1 & 0.214 & 21.46 & 37.46 \\
Renewed once & 2 & 0.239 & 21.82 & 21.82 \\
Renewed once & 3 & 0.264 & 22.19 & 22.19 \\
\bottomrule
\end{tabular}
\end{table}

The unrepaired path charges \(32.36 + 34.12 + 35.99 = 102.5\) minutes and finishes the third cycle at health \(0.676\), with failure probabilities \(0.106\), \(0.132\), and \(0.159\). The repaired path pays \(16 + 21.46 = 37.46\) minutes on the first cycle, which is slower than the unrepaired first cycle, and then \(21.82\) and \(22.19\) minutes with failure probability zero, because health stays below \(0.40\). The three-cycle total is \(81.5\) minutes. The saving is about twenty-one minutes, and the cutter that remains is still a healthy fast machine rather than a machine heading toward \(0.68\).

Figure~\ref{fig:threeh} draws the two health paths. The maintenance step is the discontinuity. Everything after it is ordinary wear at the cutter's drift of \(0.00115\) per minute. A planner who looks only at the first cycle, where maintenance is the slower option in Table~\ref{tab:worked}, will refuse the twin. A planner who looks at the third cycle will not. The estimation in the wear scenario is the third-cycle view, repeated across every machine that crosses the tolerance.

\begin{figure}[ht]
\centering
\begin{tikzpicture}
\begin{axis}[
  width=0.86\textwidth,
  height=6.2cm,
  xlabel={Shell cut},
  ylabel={Health at the start of the cut},
  xmin=0.8, xmax=3.2,
  ymin=0, ymax=0.8,
  xtick={1,2,3},
  grid=both,
  legend pos=north west,
  thick
]
\addplot coordinates {(1,0.580) (2,0.611) (3,0.643)};
\addlegendentry{No maintenance}
\addplot coordinates {(1,0.214) (2,0.239) (3,0.264)};
\addlegendentry{Renewed before the first cut}
\end{axis}
\end{tikzpicture}
\caption{Health of Cutter-A at the start of three identical shell cuts. The renewed path begins after Equation~\eqref{eq:pm-map}. Neither path is a failure time. Both are the state the next score will read.}
\label{fig:threeh}
\end{figure}

The energy accounting moves the same way and is easier to underestimate. The unrepaired first cycle draws about \(3.90\,\mathrm{kWh}\) plus a share of the repair. The repaired first cycle draws \(4.65\,\mathrm{kWh}\) because the maintenance minutes are charged. By the third cycle the unrepaired cutter is both slower and hotter, and the nameplate \(6.8\,\mathrm{kW}\) is multiplied by a larger health surcharge. The day-level energy gaps in Chapter~\ref{ch:estimation}, on the order of \(100\,\mathrm{kWh}\), are this third-cycle effect summed over every occupation, not a property of one row of Table~\ref{tab:worked}.

\subsection{How many decisions a day contains}

A book of sixteen orders has, on the routes of Section~\ref{sec:problem-set}, somewhere between six and eight operations each, so the cell starts on the order of one hundred and twenty bundles. Each bundle is one call to the scorer. At a typical call the feasible set is not one hundred and twenty. Precedence keeps only the next operation of each released order, and capacity keeps only the free machines and free operators. A call with eight ready operations, two legal machines, and two legal operators is a few dozen scores, each a constant-time evaluation of Equations~\eqref{eq:raw1}--\eqref{eq:psi} and, when the twin is on, Algorithm~\ref{alg:ssdt}. The day is a few thousand such scores. That is why identification, which repeats the day a few thousand times, finishes as a batch job, and why dispatch, which does it once, is not the thing a planner waits for.

The multi-layer model changes the constant, not the number of decisions. Each layer costs a pass over the edges of the candidate graph. Four nodes and a handful of typed edges remain a small graph. The reason to defer that model is not computer time. It is the absence, in this thesis, of a plant large enough for the hand-grouped keys to be the wrong inductive bias. Twenty identified coefficients already expose which costs matter. A deeper network would hide that exposure unless its attention weights were read back out, which is possible and was not necessary here.

The evolutionary reference pays a different bill. A population of eighteen, decoded twelve times and then doubled inside the generational pool, is a few hundred full days of simulation for one book, and it returns a front for that book only. Repeating it at every rush release is the cost Section~\ref{sec:tech1-review} objected to. The online score pays the identification once and then pays one argmin per free resource.

\subsection{Where the disjunctive graph went}

It is reasonable to ask whether the four-node candidate graph replaced the disjunctive graph or merely decorated it. Precedence is still a chain of conjunctive arcs: an operation enters \(R(t)\) only when its predecessor has finished. Machine capacity is still disjunctive: two bundles share a machine only in sequence, because a busy machine is not in \(\mathcal{F}\). The same sentence holds for an operator. The token inequality \eqref{eq:token} has no classical disjunctive counterpart, because a die is not a machine and two operations can hold two dies at the same time. The candidate graph is the disjunctive graph at the moment a decision is taken, restricted to the operations that could actually start, and extended by the operator and the shop node.

What was given up is an explicit binary \(y_{aa'}\) for every pair of operations that might ever conflict. Those binaries are the right variables for a solver that wants a certificate of optimality on a frozen instance. They are the wrong variables for a policy that must act when the instance is not frozen. Proposition~\ref{prop:mask} is the certificate this construction offers instead: every executed bundle was feasible when it started, and every operation eventually starts. It is not a certificate that no other feasible schedule has a smaller \(J_\omega\). The thesis does not claim that certificate. The tables measure the distance to four rules and to two ablations, which is a different and, for a running cell, a more relevant distance.

\subsection{A sensor reading, worked as a Kalman step}

The estimation plant writes health with Equation~\eqref{eq:plant-wear} and never opens Equation~\eqref{eq:kalman-upd}. This subsection runs the filter once, on the same cutter, so the sensor path is not only a citation. Take the posterior before a shell cut to be \(\hat h = 0.58\) with variance \(v = 0.010\). The drift is the cutter's own \(\eta = 0.00115\) per minute. The diffusion scale is taken as \(\sigma = 0.015\) per square-root minute, a design value for the indicator noise, not a fitted plant parameter. The cycle length is the unrepaired \(p = 27.268\) minutes.

Predict:
\begin{equation}
\hat h^{-} = 0.58 + 0.00115 \times 27.268 = 0.6114 ,
\end{equation}
\begin{equation}
v^{-} = 0.010 + (0.015)^2 \times 27.268 = 0.01614 .
\end{equation}
Suppose the indicator, after the feature map, reads \(y = 0.64\), a little worse than the drift expected, with noise variance \(\rho^2 = 0.008\). The gain and the update are
\begin{equation}
K = \frac{0.01614}{0.01614 + 0.008} = 0.6685 ,
\end{equation}
\begin{equation}
\hat h = 0.6114 + 0.6685 \times (0.64 - 0.6114) = 0.6305 ,
\end{equation}
\begin{equation}
v = (1 - 0.6685) \times 0.01614 = 0.00535 .
\end{equation}
The reading moved the state from the pure-drift value \(0.611\) to \(0.631\), and it cut the variance. A second reading will move it less, because \(v\) is now smaller and the gain will fall. That is the filter doing the job Gebraeel and co-authors asked of a Bayesian residual: one bad minute does not become the new truth, and a persistent bias does not stay outside the state \cite{gebraeel2005}.

Now apply the Gaussian ceiling \eqref{eq:gauss-threshold} at \(\varepsilon = 0.10\), so \(\Phi^{-1}(0.90) \approx 1.282\), with ceiling \(h^{\star} = 0.90\), to the question ``may this same cycle be repeated.'' Using the updated health,
\begin{equation}
0.6305 + 0.00115 \times 27.268 + 0.015 \sqrt{27.268} \times 1.282 = 0.762 .
\end{equation}
The value \(0.762\) is under \(0.90\). The Gaussian one-step ceiling would allow another cycle. The hazard \eqref{eq:phi} at health \(0.6305\) and the same cycle length is \(0.144\), which is above \(\varepsilon = 0.10\). Algorithm~\ref{alg:ssdt} would veto. The two forms, which Section~\ref{sec:ssdt} presented as the same kind of object, disagree on this borderline cycle.

The disagreement is useful and it is not a bug in the arithmetic. Equation~\eqref{eq:gauss-threshold} asks whether the health at the end of the cycle is likely to have crossed \(0.90\). Equation~\eqref{eq:phi} asks whether an interruption intensity, which turns on at \(0.40\), has accumulated enough hazard to exceed the tolerance. A machine can be likely to survive the cycle without hitting \(0.90\) and still be likely enough to suffer a stop that the factory does not want. The estimation uses the hazard, because the cell's interruption model is an intensity, not a first hitting time of \(0.90\). A plant whose maintenance standard is literally ``stop at indicator \(0.90\)'' should use the Gaussian ceiling instead, and should expect fewer vetoes on this cutter. Publishing both formulas, and one numerical case where they split, is what makes the tolerance a parameter rather than a story.

The variance also says when the filter is not yet allowed to veto. At the start of a shift, if \(v\) is large, \(K\) is near one and a single noisy \(y\) can throw \(\hat h\) across \(\varepsilon\). The practical rule is to withhold Algorithm~\ref{alg:ssdt} until \(v\) has fallen below a band, or to inflate \(\varepsilon\) by a term in \(\sqrt{v}\). Neither guard is in the estimation, because the plant variance is zero. Both belong in the specification for a cutter whose indicator is new.

\subsection{What Chapter~\ref{ch:estimation} is allowed to claim about these formulas}

\subsection{Failure modes of the score}

The score can be wrong in ways the mask will not catch, because the mask only enforces feasibility. Four failure modes are worth naming before the estimation, so that a bad table can be traced to a cause.

The first is a stale divisor. If batch quantities move from around eighty to around twenty, \(\rho_1 = p/45\) becomes a small number for every candidate and the processing-time weight stops discriminating. The policy then leans on whatever cost is still order-one, which on this vector is load deviation. The shop will look balanced and late. The fix is to re-identify, not to edit \(5.660\) by hand.

The second is a health indicator that saturates. If every machine reports \(h\) near one because the feature map was fit on a different tool, then \(\rho_5\) and \(\rho_8\) are large for every candidate and the maintenance test fires on every cycle. The day fills with sixteen-minute renewals that do not change a saturated indicator, because Equation~\eqref{eq:pm-map} pulls \(h\) down and the next reading, if it is still saturated, puts it back. The twin then oscillates. The fix is the feature map, upstream of the scheduler.

The third is an opening that was identified at one preference and then used at another. Section~\ref{sec:rch} already recorded that the shop-node opening at the balanced \(\omega\) is negative and nearly silent. Feeding an energy-heavy \(\omega\) through Equations~\eqref{eq:open-op}--\eqref{eq:open-s} without a new identification does not create an energy policy. It rescales a policy that was fit to balance load. Chapter~\ref{ch:estimation} therefore does not draw a Pareto front from those openings.

The fourth is a token capacity of zero that the data claim is one. The mask will refuse every operation of that colour, the clock will advance only when some other colour can run, and an order of the missing colour will wait until the end and then, if no other work remains, deadlock is avoided only because capacity one was assumed in Proposition~\ref{prop:mask}. A capacity of zero violates the hypothesis of the proposition. The cell should reject such a book at load time rather than discover it as a stuck clock.

None of these modes appeared in the estimation books, because the divisors match the batch sizes, the health values start inside \((0,1)\) and move by the wear recursion rather than by a saturated sensor, the preference is the one that was identified, and every token pool has capacity at least one. They are listed because a transfer of the same code to a different bag line will hit them in that order.

\begin{table}[ht]
\centering
\caption{What to check when a transferred cell misbehaves. The estimation books did not trip these checks.}
\label{tab:failmodes}
\small
\begin{tabular}{llp{6.2cm}}
\toprule
Symptom & Likely cause & Check \\
\midrule
Balanced and late & stale time divisor & re-identify after rescaling \(\rho_1\) \\
Maintenance every cycle & saturated health & inspect the feature map, not \(\varepsilon\) \\
Energy weight ignored & old \(\omega\) in the openings & rerun Algorithm~\ref{alg:cem} \\
Clock stuck on one colour & token capacity zero & reject the book before dispatch \\
\bottomrule
\end{tabular}
\end{table}

An implementation that follows this chapter, rather than the estimation program, has a definite order of work. First the routes, the skill matrix, the token capacities, and the nameplate powers are tables, not code. Second the mask is written and tested on a first-in-first-out ranking, because Proposition~\ref{prop:mask} can be checked before any weight exists: every order must finish, and no token count may go negative. Third the twelve raw costs are logged for a single candidate and checked against a hand calculation of the kind in the cutter example. Fourth the weights are identified, and the log-weights are inspected. A colour weight near zero, as in Equation~\eqref{eq:wstar}, is a result to be reported, not a bug to be patched by hand. Fifth the maintenance test is switched on at a stated \(\varepsilon\), and a one-machine trace like Table~\ref{tab:threecycle} is checked before any full book is compared. Sixth, and only if a sensor exists, the Kalman step is put in front of the wear recursion and the variance guard is set so that a new indicator cannot veto on its first noisy hour. Skipping to a multi-layer network at step two would make every later check harder and would not change the mask.

The two technologies, stripped of the surrounding discussion, are small enough to restate. RCH-GAT builds, at each free resource, the feasible bundles of Equation~\eqref{eq:bundle}, scores them by the preference-weighted attention of Equation~\eqref{eq:psi}, and starts the minimiser. The graph has four node types because the bundle names an operation, a machine, an operator, and a token. The weights in Equation~\eqref{eq:wstar} say which of those nodes the identification actually used: load, skill, failure, setup, and energy, with colour contention nearly silent. SSDT-CC writes the health that those failure and energy costs read. It lengthens processing time by Equation~\eqref{eq:ptime}, assigns an interruption probability by Equation~\eqref{eq:phi}, and, when that probability exceeds \(\varepsilon\) or the predicted health exceeds \(0.90\), renews the machine by Equation~\eqref{eq:pm-map} before the cycle. The renewal is not free on the cycle that pays it. Three cycles later, on the numbers of Table~\ref{tab:threecycle}, it is cheaper than letting the cutter drift. A sensor path, when one exists, revises the same health by Equations~\eqref{eq:kalman-pred} and \eqref{eq:kalman-upd}, and it can disagree with the hazard on a borderline cycle, as the numerical Kalman step showed. The estimation is allowed to speak only for the subset it ran.

The formulas of this chapter define the system. The next chapter measures an executable subset of them on a simulated school-bag cell. The subset includes Equations~\eqref{eq:ptime}, \eqref{eq:raw1}--\eqref{eq:psi}, \eqref{eq:phi}, \eqref{eq:pm-map}, and Algorithms~\ref{alg:cem}--\ref{alg:ssdt}. The subset does not include a trained multi-layer network \eqref{eq:gat-full}--\eqref{eq:gat-update} and does not include a noisy sensor around \eqref{eq:kalman-upd}. Claims about makespan, energy, and failures in the next chapter are claims about the subset. Claims about the multi-layer network and the Kalman filter are claims about the model that the subset instantiates, and they are not given extra experimental authority.

\chapter{System Build and Estimation}
\label{ch:estimation}

The formulas of Chapter~\ref{ch:methods} do not, by themselves, say whether the two technologies move the five targets. This chapter builds the smallest system that can be run, and then estimates it. Section~\ref{sec:build} describes the cell, the data, the programs, and the way RCH-GAT and SSDT-CC sit inside them. Section~\ref{sec:estimate} reports what those programs produced on four stresses of the order book, on two ablations, and on a sweep of the chance-constraint tolerance. The estimates are means and standard deviations over paired random order books. They are not measurements from a physical sewing line, and the text does not treat them as if they were.

\section{System build}
\label{sec:build}

\subsection{What was built}

The system is a discrete-event shop cell and a dispatcher that share one state. The cell holds orders, machines, operators, dies, and thread lots. The dispatcher is called whenever the cell has a free feasible bundle. It returns one bundle, optionally prefixed by preventive maintenance. The cell advances that bundle, writes the new health, releases the resources at the end, and calls the dispatcher again. There is no nightly plan file. Figure~\ref{fig:architecture} is the whole architecture.

\begin{figure}[ht]
\centering
\begin{tikzpicture}[
  box/.style={draw, rounded corners, align=center, font=\small, minimum width=2.7cm, minimum height=0.9cm},
  arr/.style={-{Stealth}, thick}
]
\node[box] (ord) {Order book\\releases, routes};
\node[box, right=1.0cm of ord] (cell) {Shop cell\\machines, people, tokens};
\node[box, right=1.0cm of cell] (twin) {Twin\\health and \(\phi\)};
\node[box, below=1.05cm of twin] (pol) {RCH-GAT\\masked score};
\node[box, below=1.05cm of cell] (log) {KPI log\\five targets};
\draw[arr] (ord) -- (cell);
\draw[arr] (cell) -- (twin);
\draw[arr] (twin) -- (pol);
\draw[arr] (pol) -- (cell);
\draw[arr] (cell) -- (log);
\end{tikzpicture}
\caption{Executable architecture. The policy reads the twin and writes one bundle back into the cell. The KPI log is the only path from a run to the tables in Section~\ref{sec:estimate}.}
\label{fig:architecture}
\end{figure}

Four programs implement the figure. The shop generator builds an order book from a seed and a scenario name. The cell executes Algorithm~\ref{alg:dispatch}. The scorer evaluates Equation~\eqref{eq:psi} and Algorithm~\ref{alg:ssdt}. The identification program runs Algorithm~\ref{alg:cem} and stores the coefficient vector. A comparison driver repeats the cell for every method, every scenario, and every seed, and writes the means. The driver is deterministic given the seed. Re-running it reproduces the tables.

The cell is the plant. It is also the sensor, in the sense of Section~\ref{sec:ssdt}: health is updated by the wear recursion \eqref{eq:plant-wear} rather than by a Kalman correction of a noisy current trace. That choice is deliberate. It removes sensor luck from the comparison. It does not remove the status effect, because wear, the duration factor, the hazard, and the maintenance map are all active.

\subsection{The cell}

Nine machines form seven groups. Table~\ref{tab:machines} lists the quantities that enter Equations~\eqref{eq:ptime} and \eqref{eq:f3}. Health at the start of the day is part of the scenario, not a warm-up artefact. Cutter-A is fast, powerful, and already at \(0.58\). Lockstitch-2 is at \(0.52\). The embroidery head is at \(0.46\). A scheduler that ignores status is ignoring the largest machines in the cell, not a rounding error on a healthy fleet.

\begin{table}[ht]
\centering
\caption{Machines in the estimation cell. Wear is the drift \(\eta_k\) in health per processing minute. Health \(h_0\) is the nominal-scenario initial state.}
\label{tab:machines}
\small
\begin{tabular}{llrrrr}
\toprule
Machine & Group & Power (kW) & Speed & Wear & \(h_0\) \\
\midrule
Cutter-A & cut & 6.8 & 1.28 & 0.00115 & 0.58 \\
Cutter-B & cut & 2.4 & 0.74 & 0.00055 & 0.22 \\
Embroidery-1 & print & 3.4 & 1.00 & 0.00050 & 0.46 \\
Lockstitch-1 & sew & 0.75 & 1.05 & 0.00038 & 0.24 \\
Lockstitch-2 & sew & 0.75 & 0.98 & 0.00040 & 0.52 \\
Bartack-1 & sew, heavy & 1.15 & 1.00 & 0.00048 & 0.34 \\
Hardware-1 & hardware & 1.45 & 1.00 & 0.00028 & 0.20 \\
QC-1 & inspect & 0.35 & 1.00 & 0.00008 & 0.08 \\
Pack-1 & pack & 0.55 & 1.00 & 0.00008 & 0.12 \\
\bottomrule
\end{tabular}
\end{table}

Eight operators staff the nine machines, which is the dual-resource constraint. Table~\ref{tab:operators} is the skill matrix. A blank entry is a hard zero and is masked. Cut-2 can pack slowly and is the only second cutter. Print-1 is the only embroidery operator and can stand in at inspection. Sew-3 cannot run the bartack. The hardware operator and the inspection operator both reach the packing table, so packing is rarely the resource that deadlocks, and hardware sometimes is.

\begin{table}[ht]
\centering
\caption{Operator skills. The number multiplies machine speed in Equation~\eqref{eq:ptime}. A dash is an infeasible pairing.}
\label{tab:operators}
\small
\begin{tabular}{lrrrrrrr}
\toprule
Operator & Cut & Print & Sew & Heavy & Hardware & Inspect & Pack \\
\midrule
Cut-1 & 1.00 & -- & -- & -- & -- & -- & -- \\
Cut-2 & 0.82 & -- & -- & -- & -- & -- & 0.55 \\
Print-1 & -- & 1.00 & -- & -- & -- & 0.45 & -- \\
Sew-1 & -- & -- & 1.06 & 0.78 & -- & -- & -- \\
Sew-2 & -- & -- & 0.96 & 1.08 & -- & -- & -- \\
Sew-3 & -- & -- & 1.00 & -- & -- & -- & -- \\
Hard-1 & -- & -- & -- & -- & 1.00 & -- & 0.70 \\
QC-1 & -- & -- & -- & -- & 0.62 & 1.00 & 0.95 \\
\bottomrule
\end{tabular}
\end{table}

Tokens are few on purpose. Two dies cover two cutters. Thread lots allow two concurrent users of navy, black, red, and grey, and one user of khaki. In the shortage scenario the navy, red, and khaki pools drop to one. Colour setup is eleven minutes on cutting, embroidery, and sewing when the previous colour differs, and four minutes when the machine has no remembered colour. The repair expectation, when a failure probability is priced, is forty-eight minutes at half nameplate power. Preventive maintenance, when the chance constraint fires, is sixteen minutes and applies Equation~\eqref{eq:pm-map}.

Routes are those of Section~\ref{sec:problem-set}. A backpack has eight operations, a trolley bag seven, and a sports sack six. Base minutes at quantity 80 and size M are stored with the route. A random order draws a style, a colour, a size in \(\{S,M,L\}\), a quantity in \(\{40,60,80,100,120\}\), a release time in \(\{0,40,90,150\}\) minutes, and a due date equal to the release time plus a multiple of the route's base work. The multiple is drawn uniformly from \(3.6\) to \(6.4\), which is wide enough that a good schedule can finish a substantial share of orders on time and a poor schedule cannot. Sixteen orders form a book, except in the rush scenario, which adds four more at time \(220\).

\subsection{Methods that share the cell}

Every method uses the same cell, the same mask, and the same health plant. They differ in the ranking of feasible bundles and in whether Algorithm~\ref{alg:ssdt} is allowed to maintain.

First-in-first-out ranks by release time, then by order identity, then by operation index. Shortest processing time ranks by processing time plus setup, using the health-scaled time of Equation~\eqref{eq:ptime}, because hiding health inside the baseline's clock and not inside its ranking would invent a gain. The baselines are allowed to see the true processing time. They are not given the failure features, the token-contention features, or the maintenance action. Earliest due date ranks by \(d_j\). Apparent tardiness cost uses Equation~\eqref{eq:atc} with look-ahead \(k = 2\) and a forty-minute average processing time.

RCH is Equation~\eqref{eq:psi} with the identified coefficients and with the maintenance test switched off. RCH-SSDT is the same score with the test switched on at \(\varepsilon = 0.10\). The no-resource ablation zeroes the contention costs \(\rho_9,\rho_{10},\rho_{11}\) and the shop-node key, and still allows maintenance. The no-status ablation zeroes health and failure in the score and switches maintenance off. The plant's wear and the plant's expected repair delay remain active for every method, including the ablations. Otherwise a method that ``ignores status'' would be running a different factory.

An evolutionary reference is reported separately and is not given maintenance. On four nominal books a small NSGA-II \cite{deb2002}, with a population of eighteen and twelve generations, searches a priority key per operation and a machine bias. The decoder is the same cell. The front it returns is a static search on one book, not an online policy. It is there to show what a short population search can move when it is allowed to look at the whole book at once, and what it does not move when it cannot renew a machine.

\subsection{Identification}

Coefficients were identified on four books, seeds \(11\), \(23\), \(37\), and \(41\), which are not in the estimation seeds. The preference during identification was the balanced vector \(\omega = (0.34, 0.33, 0.18, 0.15)\), ordered as makespan, tardiness, energy, and imbalance, with the failure ratio added at weight \(0.12\) as in Equation~\eqref{eq:scalar}. The cross-entropy search used a population of \(36\), an elite of \(8\), and \(28\) generations. Figure~\ref{fig:train} shows the best and the mean fitness. Fitness is \(J_\omega\) relative to first-in-first-out, so a value of one is a tie with that rule and a value below one is an improvement on the scalar.

\begin{figure}[ht]
\centering
\begin{tikzpicture}
\begin{axis}[
  width=0.92\textwidth,
  height=6.2cm,
  xlabel={Generation},
  ylabel={Scalar cost relative to FIFO},
  ymin=0.30, ymax=0.65,
  legend pos=north east,
  grid=both,
  every axis plot/.append style={thick}
]
\addplot coordinates {
(0,0.450) (1,0.415) (2,0.425) (3,0.400) (4,0.403) (5,0.405)
(6,0.400) (7,0.398) (8,0.397) (9,0.387) (10,0.402) (11,0.399)
(12,0.386) (13,0.395) (14,0.397) (15,0.391) (16,0.396) (17,0.388)
(18,0.390) (19,0.369) (20,0.379) (21,0.372) (22,0.373) (23,0.373)
(24,0.371) (25,0.372) (26,0.376) (27,0.366)
};
\addlegendentry{Best in generation}
\addplot coordinates {
(0,0.592) (1,0.495) (2,0.469) (3,0.456) (4,0.453) (5,0.442)
(6,0.429) (7,0.420) (8,0.423) (9,0.418) (10,0.421) (11,0.426)
(12,0.416) (13,0.419) (14,0.423) (15,0.412) (16,0.410) (17,0.415)
(18,0.414) (19,0.408) (20,0.401) (21,0.395) (22,0.397) (23,0.393)
(24,0.397) (25,0.389) (26,0.404) (27,0.389)
};
\addlegendentry{Population mean}
\end{axis}
\end{tikzpicture}
\caption{Identification of the twenty score coefficients. The vertical axis is the preference-weighted cost divided, term by term, by first-in-first-out on the same training books. The search was stopped at generation 27. The coefficient vector used in the tables is the best vector observed, at \(0.366\).}
\label{fig:train}
\end{figure}

The curve falls quickly and then flattens. Most of the scalar gain is present by generation ten. Later generations trim the best value from about \(0.39\) to \(0.366\) and do not change the identity of the method. This is the behaviour one expects when the score is a low-dimensional monotone function of sensible costs and the mask already carries the combinatorial structure. It is not evidence that a deep multi-layer network would be unnecessary on a plant with hundreds of machines. It is evidence that, on this cell, the attention score of Section~\ref{sec:rch} can be identified without hiding the result inside an opaque training run.

\subsection{Scenarios and the statistics that are reported}

Four scenarios stress the problem set.

The nominal scenario is the cell as tabulated, with sixteen orders and the standard token capacities. It is the fair day.

The rush scenario releases four extra orders at minute \(220\), with due dates at \(2.15\) times their base work after that release. They are tighter than the ordinary multiple of \(3.6\) to \(6.4\). Every online method, including the classical rules, can see them only after minute \(220\). The scenario asks whether rebuilding the feasible set, which all online methods do, is enough, or whether the score and the twin still matter once the book has grown.

The wear scenario multiplies every wear rate by \(1.85\) and adds \(0.10\) to every initial health. Cutter-A then starts at \(0.68\), already deep into the hazard. This is the scenario the twin is for.

The shortage scenario leaves wear nominal and cuts the navy, red, and khaki pools to one concurrent user. This is the scenario the shop node is for.

Each scenario is run on twenty books, seeds \(1000\) through \(1019\), paired across methods. Tables report the mean and, in the text, the standard deviation. Total tardiness has a large standard deviation because one late trolley bag dominates the sum. A difference of a few tens of minutes in mean tardiness, against a standard deviation of a few hundreds, is reported as a point estimate and is not promoted to a ranking. Differences that exceed the spread, in particular makespan, energy, and expected failures for the full system against the classical rules, are treated as the estimation's actual findings.

The failure column is an expectation, summed from Equation~\eqref{eq:phi}. Two methods can be compared on it without a second random draw. Expected repair minutes are inserted into the timeline so that a method which accepts a high \(\phi\) also accepts the blocking time that follows.

\section{Estimation}
\label{sec:estimate}

\subsection{Nominal book}

Table~\ref{tab:nominal} is the main estimate. Times are minutes. Energy is kilowatt-hours. Failures are the expected count \(f_5\). On-time rate is the fraction of orders with no tardiness. The twenty-book standard deviation of makespan is about \(166\) minutes for first-in-first-out and about \(83\) minutes for the full system. The standard deviation of tardiness is several hundred minutes for every method.

\begin{table}[ht]
\centering
\caption{Nominal order books, twenty paired seeds. Lower is better except for the on-time rate.}
\label{tab:nominal}
\small
\begin{tabular}{lrrrrrrr}
\toprule
Method & Makespan & Tardiness & On time & Energy & Imbalance & Failures & PM \\
\midrule
FIFO & 1520.7 & 2761.4 & 0.519 & 303.2 & 377.7 & 18.02 & 0 \\
SPT & 1547.6 & 512.8 & 0.816 & 268.3 & 366.0 & 15.56 & 0 \\
EDD & 1545.4 & 212.1 & 0.891 & 299.3 & 377.4 & 16.95 & 0 \\
ATC & 1631.8 & 274.7 & 0.863 & 299.1 & 393.8 & 17.44 & 0 \\
RCH & 1234.4 & 260.9 & 0.853 & 211.0 & 260.4 & 10.13 & 0 \\
RCH-SSDT & 1051.6 & 165.1 & 0.903 & 167.5 & 222.1 & 0.91 & 5.35 \\
No resource & 1050.6 & 135.4 & 0.903 & 168.2 & 218.6 & 0.89 & 5.70 \\
No status & 1260.4 & 268.5 & 0.875 & 207.7 & 263.8 & 10.22 & 0 \\
\bottomrule
\end{tabular}
\end{table}

Read the classical block first. First-in-first-out finishes about half the orders on time and accumulates \(2761\) minutes of tardiness. Shortest processing time and earliest due date cut that tardiness sharply, to \(513\) and \(212\) minutes, and they do not cut the makespan. All four rules finish the day in roughly \(1520\) to \(1632\) minutes and burn roughly \(268\) to \(303\,\mathrm{kWh}\). Expected failures sit between \(15.6\) and \(18.0\). Apparent tardiness cost does not beat earliest due date on this book. The look-ahead in Equation~\eqref{eq:atc} is a good idea on a pure tardiness problem and is not, here, a substitute for health and tokens.

RCH, which is the graph score without maintenance, is the estimate of the first technology alone. Makespan falls from \(1520.7\) to \(1234.4\) minutes, a reduction of \(286\) minutes against first-in-first-out and of about \(310\) minutes against earliest due date. Energy falls from the best classical value, the \(268.3\,\mathrm{kWh}\) of shortest processing time, to \(211.0\,\mathrm{kWh}\). Imbalance falls from about \(370\) to \(260\). Expected failures fall from about \(16\) to \(10.1\). Tardiness, at \(260.9\) minutes, is better than shortest processing time and worse, as a point estimate, than earliest due date's \(212.1\). The on-time rates are \(0.85\) and \(0.89\). Given tardiness standard deviations of \(250\) and \(356\) minutes, that ordering is not a stable ranking. The honest nominal statement about RCH is that it shortens the day, saves energy, balances the machines, and avoids some failures, while staying in the same tardiness band as the due-date rules.

RCH-SSDT is the estimate of the two technologies together. Makespan falls further, to \(1051.6\) minutes. Energy falls to \(167.5\,\mathrm{kWh}\). Expected failures fall to \(0.91\), with \(5.35\) preventive maintenances on an average book. Expected breakdown delay, which is \(865\) minutes under first-in-first-out and \(486\) minutes under RCH, is \(44\) minutes under the full system. Tardiness' point estimate, \(165\) minutes, is below earliest due date, and the on-time rate is \(0.90\). The tardiness standard deviation is still \(205\) minutes, so the due-date claim is the weak one. The makespan, energy, failure, and balance claims are not weak: the full system's makespan spread (\(\pm 83\)) and the classical makespan spread (\(\pm 150\) or more) do not cover a gap of about \(470\) minutes, and the failure gap is from the teens to under one.

The ablations split the gain. Removing status features and maintenance returns the result to the RCH band: makespan \(1260\), energy \(208\), failures \(10.2\). The twin, not the attention biases, is what drives failures from ten to under one. Removing resource features and keeping maintenance does not hurt the nominal book. The point estimates of makespan and energy stay inside a few minutes of the full system, and the tardiness point estimate is even lower. Token pools of two, on a sixteen-order nominal book, are not tight enough for the shop node to earn its keep. That is a finding about this scenario, not a reason to delete the node. The shortage and rush scenarios are where the node was supposed to matter, and they are reported next.

Figure~\ref{fig:bars} repeats the four headline means so the gap is visible without the full table.

\begin{figure}[ht]
\centering
\begin{tikzpicture}
\begin{axis}[
  width=0.96\textwidth,
  height=6.4cm,
  ybar,
  bar width=6pt,
  ylabel={Value on the nominal book},
  symbolic x coords={Makespan, Energy x5, Failures x50},
  xtick=data,
  ymin=0,
  legend style={at={(0.5,-0.22)}, anchor=north, legend columns=4},
  nodes near coords style={font=\tiny},
  every node near coord/.append style={font=\tiny}
]
\addplot coordinates {(Makespan,1521) (Energy x5,1516) (Failures x50,901)};
\addplot coordinates {(Makespan,1545) (Energy x5,1497) (Failures x50,848)};
\addplot coordinates {(Makespan,1234) (Energy x5,1055) (Failures x50,506)};
\addplot coordinates {(Makespan,1052) (Energy x5,838) (Failures x50,46)};
\legend{FIFO, EDD, RCH, RCH-SSDT}
\end{axis}
\end{tikzpicture}
\caption{Nominal means rescaled onto one axis. Energy is multiplied by five and expected failures by fifty, so that a kilowatt-hour and a failure can sit beside a minute of makespan. The scaling is only for the drawing. Table~\ref{tab:nominal} has the unscaled numbers.}
\label{fig:bars}
\end{figure}

\subsection{Rush release}

Table~\ref{tab:rush} adds four orders at minute \(220\). The day is longer for every method, and tardiness grows because the new due dates are tight and because the embroidery head and the hardware bench were not idle.

\begin{table}[ht]
\centering
\caption{Rush release at minute 220, twenty paired seeds. The book contains twenty orders after the release.}
\label{tab:rush}
\small
\begin{tabular}{lrrrrrr}
\toprule
Method & Makespan & Tardiness & On time & Energy & Failures & PM \\
\midrule
FIFO & 2013.0 & 7701.6 & 0.415 & 416.9 & 30.24 & 0 \\
SPT & 1928.5 & 3090.4 & 0.515 & 356.0 & 24.30 & 0 \\
EDD & 2001.5 & 2716.0 & 0.425 & 408.2 & 27.39 & 0 \\
ATC & 2008.1 & 2871.0 & 0.410 & 403.1 & 27.11 & 0 \\
RCH & 1676.1 & 1577.6 & 0.495 & 302.4 & 17.93 & 0 \\
RCH-SSDT & 1264.6 & 735.5 & 0.685 & 214.6 & 1.02 & 7.85 \\
No resource & 1274.9 & 805.8 & 0.708 & 213.0 & 1.02 & 7.70 \\
No status & 1597.5 & 1473.6 & 0.557 & 298.0 & 17.15 & 0 \\
\bottomrule
\end{tabular}
\end{table}

Earliest due date, which won the nominal tardiness point estimate, now records \(2716\) minutes of tardiness and an on-time rate of \(0.43\). The rule does see the rush orders. Seeing them is not the same as placing them onto a machine that will still be healthy when they arrive. Shortest processing time is similar, with a slightly shorter day and slightly less energy, and with tardiness still near \(3100\) minutes. Standard deviations of rush tardiness are about \(900\) to \(1300\) minutes for the classical rules and \(522\) minutes for the full system. The gap from \(2716\) to \(736\) is large beside those spreads. The on-time rate moves from \(0.43\) to \(0.69\).

The full system finishes in \(1265\) minutes instead of about \(2000\), at \(215\,\mathrm{kWh}\) instead of about \(400\), with one expected failure instead of about twenty-seven, using \(7.9\) maintenances. RCH without maintenance already reaches \(1676\) minutes, \(302\,\mathrm{kWh}\), and \(17.9\) expected failures, so the graph score is doing part of the work before the twin renews anything. The no-status ablation sits next to that RCH result. The no-resource ablation sits next to the full system and is a little worse on makespan and on the tardiness point estimate (\(1275\) and \(806\) against \(1265\) and \(736\)). The resource features point the right way under a rush, and the difference is much smaller than the difference made by maintenance. A rush overloads machines and health before it overloads token pools.

\subsection{Accelerated wear}

Table~\ref{tab:wear} is the twin's scenario. Initial health is higher and every processing minute writes more wear.

\begin{table}[ht]
\centering
\caption{Accelerated wear, twenty paired seeds. Initial health is shifted up by \(0.10\) and every wear coefficient is multiplied by \(1.85\).}
\label{tab:wear}
\small
\begin{tabular}{lrrrrrr}
\toprule
Method & Makespan & Tardiness & On time & Energy & Failures & PM \\
\midrule
FIFO & 2097.7 & 5898.7 & 0.397 & 388.2 & 42.26 & 0 \\
SPT & 2034.0 & 2093.5 & 0.641 & 351.7 & 36.24 & 0 \\
EDD & 2059.3 & 1849.0 & 0.584 & 384.3 & 38.75 & 0 \\
ATC & 2120.8 & 2252.4 & 0.547 & 382.0 & 38.91 & 0 \\
RCH & 1782.0 & 895.1 & 0.697 & 310.0 & 30.58 & 0 \\
RCH-SSDT & 1121.7 & 287.2 & 0.850 & 176.1 & 0.81 & 11.95 \\
No resource & 1109.8 & 279.5 & 0.831 & 174.6 & 0.83 & 11.85 \\
No status & 1676.9 & 1054.7 & 0.706 & 283.6 & 28.88 & 0 \\
\bottomrule
\end{tabular}
\end{table}

Without maintenance the day runs away. Classical makespans sit near \(2100\) minutes, with \(36\) to \(42\) expected failures and about \(1700\) to \(2000\) minutes of expected repair delay charged into the timeline. RCH, which can see health and still cannot renew it, improves the book to \(1782\) minutes and \(30.6\) expected failures. Seeing the worn cutter and preferring the healthy one is worth several hundred minutes. It is not worth what renewal is worth.

With the chance constraint on, the same books finish in \(1122\) minutes, at \(176\,\mathrm{kWh}\), with \(0.81\) expected failures and \(12.0\) maintenances. The no-status ablation, at \(1677\) minutes and \(28.9\) expected failures, confirms that the score's health features without the maintenance action do not reproduce this result. The no-resource ablation is statistically the same as the full system on this scenario. Wear does not act through the khaki lot.

Twelve maintenances of sixteen minutes are \(192\) minutes of tool-change time. The expected repair delay avoided, relative to earliest due date, is about \(1860 - 39 = 1821\) minutes. The arithmetic of Table~\ref{tab:worked} was that one maintenance does not pay for itself inside one cycle. The arithmetic of Table~\ref{tab:wear} is that the same action, applied when the test fires, pays for itself across the day by a wide margin. Both statements are true. They refer to different horizons, which is why the twin is a day-level technology even though its test is a cycle-level test.

Tardiness follows the makespan. The full system's point estimate is \(287\) minutes and its on-time rate is \(0.85\), against \(1849\) minutes and \(0.58\) for earliest due date. The tardiness standard deviation is \(364\) minutes for the full system and \(1009\) for earliest due date. The gap remains large beside those spreads. This is the scenario in which the due-date claim is actually supported, because the failures that the due-date rule accepts are what make the orders late.

\subsection{Thread-lot shortage}

Table~\ref{tab:shortage} restores nominal wear and tightens three colour pools to a single concurrent user.

\begin{table}[ht]
\centering
\caption{Thread-lot shortage, twenty paired seeds. Navy, red, and khaki allow only one concurrent user.}
\label{tab:shortage}
\small
\begin{tabular}{lrrrrrr}
\toprule
Method & Makespan & Tardiness & On time & Energy & Failures & PM \\
\midrule
FIFO & 1456.0 & 2329.7 & 0.544 & 295.5 & 15.85 & 0 \\
SPT & 1462.0 & 350.7 & 0.850 & 260.0 & 13.47 & 0 \\
EDD & 1515.9 & 162.6 & 0.925 & 290.8 & 14.74 & 0 \\
ATC & 1564.0 & 219.2 & 0.916 & 289.2 & 15.01 & 0 \\
RCH & 1225.4 & 171.4 & 0.887 & 208.2 & 9.29 & 0 \\
RCH-SSDT & 1038.7 & 106.3 & 0.925 & 163.5 & 0.76 & 5.15 \\
No resource & 1074.6 & 113.0 & 0.938 & 165.4 & 0.79 & 5.30 \\
No status & 1204.6 & 220.6 & 0.891 & 203.2 & 9.01 & 0 \\
\bottomrule
\end{tabular}
\end{table}

The pattern of the nominal book returns, with one extra detail. The full system is again shorter, cooler, and almost free of expected failures: \(1039\) minutes, \(164\,\mathrm{kWh}\), \(0.76\) failures, on-time rate \(0.93\). Earliest due date matches the on-time rate and loses more than \(470\) minutes of makespan and about \(127\,\mathrm{kWh}\). RCH without maintenance sits in between, as before.

The resource ablation is the detail. Without contention features the makespan point estimate rises from \(1038.7\) to \(1074.6\) and tardiness from \(106.3\) to \(113.0\). The standard deviations are about \(90\) minutes of makespan and about \(150\) minutes of tardiness, so a \(36\)-minute makespan gap is a shift of the mean, not a separation that a single book would reliably show. It is, however, in the direction the shop node was built for, and it appears under shortage and, with a similar size, under rush, while it does not appear under nominal tokens or under pure wear. The estimation therefore supports a limited claim. Resource features are a small, correctly signed correction when lots are tight, and they are not the source of the large gains in makespan, energy, and failures. Those gains survive the ablation. A factory that already has ample dies and thread, and that is trying to justify the shop node alone, will not see them in a one-day book of this size. A factory whose khaki lot is a single spool will see a modest tightening of the day, on top of the much larger tightening that comes from not running the cutter into a repair.

\subsection{The tolerance as a knob}

The wear scenario was repeated on eight further books for five values of \(\varepsilon\). These seeds are not the twenty books of Table~\ref{tab:wear}, so the row \(\varepsilon = 0.10\) does not reproduce that table's mean exactly. The sweep is a curve, not a second copy of the main experiment. Table~\ref{tab:eps} and Figure~\ref{fig:eps} report it.

\begin{table}[ht]
\centering
\caption{Chance-constraint tolerance on eight accelerated-wear books. The policy is RCH-SSDT. Maintenance count is the mean number of times Algorithm~\ref{alg:ssdt} fires.}
\label{tab:eps}
\small
\begin{tabular}{rrrrr}
\toprule
\(\varepsilon\) & Makespan & Failures & Energy & Maintenances \\
\midrule
0.02 & 1145.3 & 0.046 & 182.0 & 15.00 \\
0.05 & 1151.3 & 0.238 & 182.3 & 13.88 \\
0.10 & 1170.2 & 0.813 & 195.7 & 12.63 \\
0.20 & 1187.2 & 2.741 & 198.2 & 10.38 \\
0.50 & 1335.7 & 9.773 & 228.6 & 7.00 \\
\bottomrule
\end{tabular}
\end{table}

\begin{figure}[ht]
\centering
\begin{tikzpicture}
\begin{axis}[
  width=0.78\textwidth,
  height=6.0cm,
  xlabel={Tolerance \(\varepsilon\)},
  ylabel={Expected failures},
  xmin=0, xmax=0.55,
  ymin=0, ymax=11,
  grid=both,
  thick
]
\addplot coordinates {(0.02,0.046) (0.05,0.238) (0.10,0.813) (0.20,2.741) (0.50,9.773)};
\end{axis}
\end{tikzpicture}
\caption{Expected failures against the chance-constraint tolerance on the eight-book wear sweep. Loosening the veto lets the worn machines run and the failure count returns.}
\label{fig:eps}
\end{figure}

Tightening \(\varepsilon\) from \(0.50\) to \(0.02\) cuts expected failures from \(9.77\) to \(0.05\) and raises the maintenance count from \(7.0\) to \(15.0\). Makespan moves from \(1336\) minutes down to \(1145\), not up. Energy moves with it, from \(228.6\,\mathrm{kWh}\) down to \(182.0\,\mathrm{kWh}\). On this wear parameterisation, a sixteen-minute renewal is cheap relative to the repair delay and the slowed cycles it prevents, all the way to the tightest tolerance tested. The knob behaves as a knob. It does not, inside the tested interval, produce the rising makespan that a much longer maintenance, or a much healthier fleet, would produce. That qualification belongs in the result. The technology exposes \(\varepsilon\) so that a plant with different renewal costs can set it from its own curve rather than from this one.

\subsection{A static evolutionary reference}

On four nominal books the small NSGA-II, without maintenance, produced nondominated sets of size \(14\), \(16\), \(7\), and \(18\). The best makespan in those four sets was \(1002\), \(1274\), \(1081\), and \(1219\) minutes. The best tardiness was \(649\), \(131\), \(571\), and \(326\) minutes. The best energy was \(204\), \(263\), \(237\), and \(238\,\mathrm{kWh}\). These are four separate single-book searches. A point that holds the best makespan of a front is not the point that holds the best tardiness of that front.

The reference is useful as a ceiling check and as a warning against a sloppy comparison. On one book, a population that can see every order at time zero can reach a makespan near \(1000\) minutes without any maintenance, which is in the neighbourhood of the full system's twenty-book mean of \(1052\). The energy attached to the evolutionary fronts does not enter the neighbourhood of \(168\,\mathrm{kWh}\), and the search has no failure action to bring it there. Giving the evolutionary method the maintenance map and a long enough population would be a different experiment. It was not run. The online system does not need that population at dispatch time, which is the operational distinction drawn in Section~\ref{sec:tech1-review} and left intact by these four fronts.

\subsection{Where the gain comes from}

The four scenarios support a single decomposition.

The classical rules are not close to each other, and the estimation does not replace them with a vague ``heuristic baseline.'' Earliest due date is the tardiness rule. Shortest processing time is the energy rule among the four, because it happens to favour shorter occupations on lighter machines more often than earliest due date does. First-in-first-out is the rule that shows the cost of having no score at all. Apparent tardiness cost does not dominate earliest due date once health-scaled processing times and colour setups are in the cell.

The graph score without maintenance, RCH, is responsible for a first cut of roughly \(250\) to \(350\) minutes of makespan and roughly \(50\) to \(100\,\mathrm{kWh}\), depending on the scenario, together with a partial reduction in expected failures. It does this while remaining competitive with earliest due date on nominal tardiness and clearly better than earliest due date on rush and wear tardiness. That is the estimate of resource-coupled attention on this cell. The resource half of the name is the smaller half. The attention is mostly using health, energy, slack, and processing time. Contention features add a modest, correctly signed shift when lots are short.

The twin is responsible for the second cut: failures from about ten, or about thirty in the wear scenario, down to about one, and a further several hundred minutes of makespan. The mechanism matches Table~\ref{tab:worked} extended over the day. Maintenance is a local cost. Avoided repair delay and restored machine speed are a global saving. The no-status ablation shows that merely penalising health inside \(\psi\), without the authority to renew the machine, captures only the first cut.

Preference conditioning was not re-identified at the corners of the simplex, and a gate sweep at the balanced coefficients did not move energy and tardiness in clean opposite directions. The estimation does not claim a traced Pareto front from those gates. What it claims is a scalarised policy at a stated \(\omega\), plus a tolerance \(\varepsilon\) that does trace a front between failures and maintenance count. A plant that wants a different compromise between makespan and energy can re-run Algorithm~\ref{alg:cem} at a new \(\omega\). The identification is short on this cell. Pretending that the present gates already provide that compromise would not match the runs.

\subsection{Spreads, and which gaps survive them}

Means without spreads invite a false ranking. Table~\ref{tab:nominal-std} gives the twenty-book standard deviations for the nominal scenario. The same qualitative pattern holds in the other three scenarios: makespan and energy under the full system are tighter as well as smaller, and tardiness is wide for everyone.

\begin{table}[ht]
\centering
\caption{Standard deviations on the twenty nominal books. Compare with the means in Table~\ref{tab:nominal}.}
\label{tab:nominal-std}
\small
\begin{tabular}{lrrrr}
\toprule
Method & Makespan & Tardiness & Energy & Failures \\
\midrule
FIFO & 165.5 & 703.0 & 35.0 & 3.29 \\
SPT & 229.6 & 571.9 & 40.8 & 3.81 \\
EDD & 175.2 & 355.5 & 39.7 & 3.48 \\
ATC & 188.2 & 410.7 & 37.6 & 3.29 \\
RCH & 146.3 & 250.5 & 28.6 & 2.25 \\
RCH-SSDT & 83.4 & 204.5 & 19.4 & 0.15 \\
No resource & 106.8 & 139.2 & 21.1 & 0.15 \\
No status & 179.2 & 282.8 & 24.3 & 2.00 \\
\bottomrule
\end{tabular}
\end{table}

A gap survives the spread when it is large next to both standard deviations, not when a \(t\)-test has been performed. No paired test is reported, and one should not be invented from the two margins. The comparisons that survive on that rough standard are these.

Makespan of RCH-SSDT against any classical rule, in every scenario. The nominal gap is about \(470\) minutes against a classical spread of \(150\) to \(230\) and a full-system spread of \(83\). The wear gap is about \(900\) minutes. The rush gap is about \(700\) minutes.

Energy of RCH-SSDT against any classical rule, in every scenario. The nominal gap against the best classical energy, shortest processing time, is about \(100\,\mathrm{kWh}\), against spreads of \(41\) and \(19\).

Expected failures of RCH-SSDT against every method that is not allowed to maintain. The nominal drop is from the teens to \(0.91\), with a full-system spread of \(0.15\). This is the cleanest column in the thesis.

Imbalance of RCH and of RCH-SSDT against the classical rules. The nominal drop is from about \(370\) to \(222\) or \(260\), and the full-system spread is \(27\) minutes of standard deviation. The load weight \(5.660\) in Equation~\eqref{eq:wstar} is the coefficient that produces this column.

RCH without maintenance against first-in-first-out on makespan and energy. The makespan gap of \(286\) minutes is larger than either spread, though less dramatic than the gap that appears once maintenance is allowed.

The comparisons that do not survive are equally part of the estimate. Nominal tardiness of RCH-SSDT against earliest due date is \(165\) against \(212\), with spreads \(205\) and \(356\). The point estimate favours the full system. A single unlucky book can reverse it. Nominal on-time rate, \(0.90\) against \(0.89\), is a tie. The resource ablation against the full system on the nominal book is a tie, and on the shortage book it is a shift of about \(36\) minutes of makespan against spreads near \(90\). Calling that shift a proved saving would over-read twenty books. Calling it noise would ignore that the sign is the same under rush and under shortage and is absent under wear, which is the pattern a real but small token effect would leave.

The full system is also more repeatable. Its makespan spread is about half the classical spread. Part of that tightness is mechanical: preventive maintenance removes the long random-looking repair delays that inflate some books and not others, even though those delays are expectations rather than draws. A method that sometimes accepts eight expected failures and sometimes accepts twenty-five will have a wide makespan. A method that holds the expected count near one will not. Repeatability is a target the problem statement did not list, and the tables give it anyway.

\subsection{What the weights say about the tables}

Equation~\eqref{eq:wstar} and Table~\ref{tab:weights} were identified before any estimation book was generated. They predict the tables more specifically than the architecture diagram does.

The largest weight is load deviation. Every scenario shows a lower imbalance for RCH than for the classical rules, and the full system lowers it further because renewed machines can take work that a worn machine would otherwise hoard and then lose to a repair. The second-largest relevant weights, skill gap and failure probability, are why a machine-only reading of the graph is the wrong reading. Skill never enters the classical rules except as a longer processing time after a weak pairing has already been chosen. Failure never enters them at all.

The smallest weight is colour contention. The shortage scenario was built to make that weight visible, and the visibility is modest. This is consistent rather than contradictory. A weight of \(0.133\) cannot move a day by the same amount as a weight of \(5.660\) or as a sixteen-minute renewal that deletes a forty-eight-minute expected repair. Die contention, at \(0.917\), is large enough to matter when both cutters want a die, and the nominal cell has two dies for two cutters, so the feature is often slack. A cell with one die would be a different estimate.

Setup weight \(1.029\) shows up as fewer colour changes: \(57.8\) under first-in-first-out and \(57.6\) under earliest due date, against \(43.2\) under the full system on the nominal book. Earliest due date groups nothing by colour. The score does, unless a due date or a health veto says otherwise. The reduction is about fifteen setups, each eleven minutes, which is \(165\) minutes of setup time returned to the day. That is a real slice of the makespan gap and it is not the whole gap. The rest is routing toward healthier, less loaded machines, and, for the full system, the repair time that is no longer charged.

\subsection{What a factory would instrument first}

The estimate implies an order of instrumentation, and the order is not the order in which the technologies were named.

The first instrument is a health indicator on the automatic cutter, the embroidery head, and the more worn lockstitch, mapped onto the unit interval that Equation~\eqref{eq:ptime} expects. Those three machines dominate power or dominate the route. A health number that cannot veto a cycle is a dashboard. The veto is Algorithm~\ref{alg:ssdt}, and it needs a maintenance action the operator is allowed to start without a separate work order. Sixteen minutes is the action this cell used. A plant whose knife change takes forty minutes should rerun the tolerance sweep before copying \(\varepsilon = 0.10\).

The second instrument is the operator skill table and a live free-time for each person. The skill weight is large, and the mask is what keeps Sew-3 off the bartack. A scheduler that knows machines and not people will reproduce the hidden wait described in Section~\ref{sec:rch}.

The third instrument is the die pool and the thread lots. They are cheap to count and they are already in the graph. On a one-day book of this size they are not where the hours are. They become first-priority only if a colour is routinely single-sourced and the due dates on that colour are the ones the planner is missing. The shortage table is the test: if a factory's own version of that table shows a large gap between the full score and the no-resource score, the lots are the constraint. If the gap looks like the thirty-minute shift reported here, the lots are not the constraint.

The network of Equations~\eqref{eq:gat-full} and \eqref{eq:gat-update} is not on this list. Nothing in the twenty-book cell is large enough to need a multi-layer embedding. The list is health with authority to maintain, then people, then tokens, and a deeper network only when the identification of twenty coefficients stops being enough. That is a smaller system than a generic Industry~4.0 diagram, and it is the system the numbers support.

\subsection{A note on the clock}

All times in the tables are minutes from the release of the first order. A nominal makespan of \(1052\) minutes is about \(17.5\) hours, longer than a single staffed shift. The cell as parameterised is a congested day-and-evening load: one embroidery head, one hardware bench, and a worn fast cutter. The absolute clock is a property of those capacities and of the batch sizes, not a claim that a school-bag plant cannot finish sixteen orders in a shift. A plant with two embroidery heads would show a shorter clock and, if the cutter were healthy, a smaller maintenance effect. The ratios are what the estimate is for. The full system is about thirty percent shorter than first-in-first-out on the nominal book (\(1521\) to \(1052\)) and about forty-five percent shorter on the wear book (\(2098\) to \(1122\)). Those ratios are the quantities a different capacity would be expected to preserve in direction, not in minutes.

Energy should be read the same way. One hundred and sixty-eight kilowatt-hours for the nominal full system is the cutter, the embroidery head, and the sewing machines over that long clock, including setups and the few maintenances. It is not a utility bill for a named factory. It is the integral of the nameplate powers in Table~\ref{tab:machines} under the wear surcharge. Moving work off Cutter-A is the main reason the integral falls. A plant whose automatic cutter is efficient when worn would see a smaller energy gap and would still see the failure gap, because the hazard does not depend on the nameplate.

\subsection{The four questions, answered in one table}

Section~\ref{sec:problem-set} asked four questions. Table~\ref{tab:summary} answers them as ratios, so the four scenarios can sit on one page. The ratio is the full system divided by the named baseline. A ratio below one is an improvement.

\begin{table}[ht]
\centering
\caption{RCH-SSDT divided by a baseline, four scenarios. Makespan and failures are against first-in-first-out. Energy is against shortest processing time, the best classical energy. Tardiness is against earliest due date, the best classical due-date rule. Ratios use the means of Tables~\ref{tab:nominal}--\ref{tab:shortage}.}
\label{tab:summary}
\small
\begin{tabular}{lrrrr}
\toprule
Scenario & Makespan / FIFO & Tardiness / EDD & Energy / SPT & Failures / FIFO \\
\midrule
Nominal & 0.69 & 0.78 & 0.62 & 0.05 \\
Rush & 0.63 & 0.27 & 0.60 & 0.03 \\
Wear & 0.53 & 0.16 & 0.50 & 0.02 \\
Shortage & 0.71 & 0.65 & 0.63 & 0.05 \\
\bottomrule
\end{tabular}
\end{table}

The first question was whether the coupled system could cut makespan, energy, and failures without giving the due-date rule a large tardiness advantage. On the nominal book the makespan ratio is \(0.69\), the energy ratio is \(0.62\), and the failure ratio is \(0.05\). The tardiness ratio against earliest due date is \(0.78\) as a ratio of means, and Section~\ref{sec:estimate} has already said that this particular ratio sits inside a wide spread. The system does not buy its energy by missing the promises the due-date rule was keeping. It also does not earn a precise claim to have beaten that rule on tardiness.

The second question was whether the same system holds under stress. The wear row is the strongest answer: makespan ratio \(0.53\), tardiness ratio \(0.16\), energy ratio \(0.50\), failure ratio \(0.02\). The rush row is the next strongest, especially on tardiness, where earliest due date sees the new orders and still finishes them late because the machines it prefers are the ones accumulating repair delay. The shortage row looks like the nominal row. Tight thread lots change the feasible set. They do not change the ranking of the technologies.

The third question was which technology produces which gain. The ratios in Table~\ref{tab:summary} are the coupled system. The ablations, which do not fit in the ratio table, split it. Removing maintenance returns failures from about one to about ten on the nominal book and to about thirty on the wear book, and it gives back a few hundred minutes of makespan. Removing resource features does not. The first technology, as identified, is a load-aware, skill-aware, health-aware score. The second technology is the action that renews health. Both are required for the failure column. Only the first is required for a large part of the imbalance column.

The fourth question was whether \(\varepsilon\) is a real knob. Table~\ref{tab:eps} is the answer, and it is monotone in the direction the chance constraint predicts: a tighter tolerance, more maintenances, fewer expected failures, and, on this wear parameterisation, a shorter rather than a longer day. A plant with a slower tool change would be expected to find a bottom in that curve. This plant, inside the tested range, did not.

A fifth question was left soft on purpose. Tardiness ratios in the nominal and shortage rows should be read as ``no damage to the due-date rule, inside a large book-to-book spread.'' Tardiness ratios in the rush and wear rows are large enough to survive that spread. The soft question is therefore not one answer. It depends on whether the day is being lost to wear and rush, or only to ordinary congestion.

\subsection{Limits of the estimate}

The cell is a simulator with a known wear recursion. A physical embroidery head has a health indicator that must be built from thread-break counts, temperature, and current, with the measurement path reviewed by Lei and co-authors \cite{lei2018}, and it has a Kalman residual that this estimate sets to zero. If that residual is large, the veto will chatter, and the sixteen-minute maintenance will be applied to noise. The formula for dealing with the residual is Equation~\eqref{eq:kalman-upd}. It was not exercised.

The order books contain sixteen or twenty orders and nine machines. Graph policies are attractive in the literature because they transfer to larger books \cite{zhang2020,song2023}. This estimate does not test that transfer. A coefficient vector identified on sixteen-order books is not, on the evidence here, a policy for a plant with two hundred machines. The multi-layer model of Equations~\eqref{eq:gat-full} and \eqref{eq:gat-update} is the form proposed for that plant, and it is untested in these tables.

Expected failures are not random failures. Inserting the mean repair delay makes methods comparable and removes the variance of a single Bernoulli draw. It also makes the timeline optimistic about clustering. Two failures that the expectation treats as independent half-repairs can, in a real morning, land on the same operator and block the bartack for a full repair plus the queue behind it. A stochastic replication would widen the spreads. It would not, unless the intensity model is badly wrong, reverse a gap between one expected failure and thirty.

Tardiness is a noisy sum. The nominal comparison with earliest due date should be read as ``no harm, within a large spread,'' not as a precise percentage. The wear and rush comparisons are the ones in which tardiness actually moves.

Resource coupling is real in the model and small in the estimate. That is the right outcome to publish. Inflating it to match the size of the maintenance effect would mislead a factory into instrumenting thread lots before instrumenting cutter wear. The instrumentation order implied by these tables is the other way around: health and the authority to act on health first, token contention second, and a deeper network only when the cell is larger than the one these coefficients were fit to.

The evolutionary reference is too small to be a competitor. It is a sanity check. A serious NSGA-III run of the kind Gong and co-authors used for five objectives \cite{gong2019}, with maintenance included in the chromosome, remains an open comparison and would be the right offline partner for the online policy.

The questions in Section~\ref{sec:problem-set} can be answered in one further sentence each, now that the ratios and the spreads are both on the page. The coupled system does cut makespan, energy, imbalance, and expected failures relative to the four rules, and on a nominal book it does not hand those rules a tardiness advantage that survives the spread. Under rush and under accelerated wear the same system holds, and tardiness moves with the makespan because the lateness was being produced by repair delay. The ablation that removes maintenance gives back the failure column and a large part of the makespan; the ablation that removes resource features does not. The tolerance \(\varepsilon\) trades maintenances against expected failures in the predicted direction. Those four sentences are the estimate. Everything else in the section is the condition under which they were measured.

Table~\ref{tab:seeds} fixes the seeds so that a rerun is a comparison and not a new experiment. Changing a seed changes the order book. Changing a parameter changes the cell. Either change is allowed, and neither is what the tables already report.

\begin{table}[ht]
\centering
\caption{Seeds used by the driver. Estimation seeds are paired across every method.}
\label{tab:seeds}
\small
\begin{tabular}{ll}
\toprule
Role & Seeds \\
\midrule
Identification books & 11, 23, 37, 41 \\
Nominal, rush, wear, shortage & 1000 to 1019 \\
Tolerance sweep, wear only & 4000 to 4007 \\
Evolutionary reference & 5000 to 5003 \\
\bottomrule
\end{tabular}
\end{table}

The runs are reproducible from the thesis directory. The cell, the score, the identification, and the comparison driver are the program \texttt{simulation/shop\_sim.py}. Training books are seeds \(11\), \(23\), \(37\), and \(41\). Estimation books are seeds \(1000\) through \(1019\). The tolerance sweep uses seeds \(4000\) through \(4007\) on the wear scenario. The evolutionary reference uses seeds \(5000\) through \(5003\). The coefficient vector is written beside the tables in \texttt{data/results.json}. A reader who changes \(\varepsilon\), a wear coefficient, or a token capacity and reruns the driver will get a different table, and that is the intended use of the program. The numbers in this chapter are the numbers that driver produced for the parameterisation written down in Section~\ref{sec:build}.

Nothing in the tables comes from a school that can be visited. The routes, the powers, and the wear rates are a constructed cell, documented so that the runs can be repeated. The technologies are estimated on that cell. The estimate supports the design: a masked resource-bundle score shortens and cools the day, and a chance-constrained renewal of machine health is what keeps the day from being spent inside expected repairs. It does not support a claim about a particular factory's next quarter.

What the cell does support is a decision about where to spend the next engineering month. Put a health indicator and a sixteen-minute maintenance authority on the worn, high-power machines before building a deeper network and before instrumenting every thread lot. Keep the operator skill table in the mask, because an infeasible bartack assignment is a hidden wait. Treat colour contention as a feature that is present and, on this book, nearly inactive, and re-identify if a future book shows the shortage gap growing past the spreads in Table~\ref{tab:nominal-std}. Leave \(\varepsilon\) visible. The wear sweep is the evidence that the number matters, and it is also the evidence that the right number depends on how long a renewal takes. A scheduler that hides the tolerance inside a trained weight cannot be retuned when the knife-change time changes. The system in this thesis can.

\begin{thebibliography}{99}

\bibitem{garey1976}
M.~R. Garey, D.~S. Johnson, and R.~Sethi.
The complexity of flowshop and jobshop scheduling.
\emph{Mathematics of Operations Research}, 1(2):117--129, 1976.

\bibitem{manne1960}
A.~S. Manne.
On the job-shop scheduling problem.
\emph{Operations Research}, 8(2):219--223, 1960.

\bibitem{giffler1960}
B.~Giffler and G.~L. Thompson.
Algorithms for solving production-scheduling problems.
\emph{Operations Research}, 8(4):487--503, 1960.

\bibitem{roy1964}
B.~Roy and B.~Sussmann.
Les probl\`emes d'ordonnancement avec contraintes disjonctives.
Note DS no.~9 bis, SEMA, Paris, 1964.

\bibitem{adams1988}
J.~Adams, E.~Balas, and D.~Zawack.
The shifting bottleneck procedure for job shop scheduling.
\emph{Management Science}, 34(3):391--401, 1988.

\bibitem{jain1999}
A.~S. Jain and S.~Meeran.
Deterministic job-shop scheduling: Past, present and future.
\emph{European Journal of Operational Research}, 113(2):390--434, 1999.

\bibitem{pinedo2016}
M.~L. Pinedo.
\emph{Scheduling: Theory, Algorithms, and Systems}.
Springer, 5th edition, 2016.

\bibitem{brucker1990}
P.~Brucker and R.~Schlie.
Job-shop scheduling with multi-purpose machines.
\emph{Computing}, 45(4):369--375, 1990.

\bibitem{brandimarte1993}
P.~Brandimarte.
Routing and scheduling in a flexible job shop by tabu search.
\emph{Annals of Operations Research}, 41(3):157--183, 1993.

\bibitem{kacem2002}
I.~Kacem, S.~Hammadi, and P.~Borne.
Approach by localization and multiobjective evolutionary optimization for flexible job-shop scheduling problems.
\emph{IEEE Transactions on Systems, Man, and Cybernetics, Part C}, 32(1):1--13, 2002.

\bibitem{pezzella2008}
F.~Pezzella, G.~Morganti, and G.~Ciaschetti.
A genetic algorithm for the flexible job-shop scheduling problem.
\emph{Computers and Operations Research}, 35(10):3202--3212, 2008.

\bibitem{gao2008}
J.~Gao, L.~Sun, and M.~Gen.
A hybrid genetic and variable neighborhood descent algorithm for flexible job shop scheduling problems.
\emph{Computers and Operations Research}, 35(9):2892--2907, 2008.

\bibitem{li2016}
X.~Li and L.~Gao.
An effective hybrid genetic algorithm and tabu search for flexible job shop scheduling problem.
\emph{International Journal of Production Economics}, 174:93--110, 2016.

\bibitem{chaudhry2016}
I.~A. Chaudhry and A.~A. Khan.
A research survey: review of flexible job shop scheduling techniques.
\emph{International Transactions in Operational Research}, 23(3):551--591, 2016.

\bibitem{panwalkar1977}
S.~S. Panwalkar and W.~Iskander.
A survey of scheduling rules.
\emph{Operations Research}, 25(1):45--61, 1977.

\bibitem{blackstone1982}
J.~H. Blackstone, D.~T. Phillips, and G.~L. Hogg.
A state-of-the-art survey of dispatching rules for manufacturing job shop operations.
\emph{International Journal of Production Research}, 20(1):27--45, 1982.

\bibitem{vepsalainen1987}
A.~P.~J. Vepsalainen and T.~E. Morton.
Priority rules for job shops with weighted tardiness costs.
\emph{Management Science}, 33(8):1035--1047, 1987.

\bibitem{deb2002}
K.~Deb, A.~Pratap, S.~Agarwal, and T.~Meyarivan.
A fast and elitist multiobjective genetic algorithm: NSGA-II.
\emph{IEEE Transactions on Evolutionary Computation}, 6(2):182--197, 2002.

\bibitem{zitzler1999}
E.~Zitzler and L.~Thiele.
Multiobjective evolutionary algorithms: A comparative case study and the strength {P}areto approach.
\emph{IEEE Transactions on Evolutionary Computation}, 3(4):257--271, 1999.

\bibitem{zitzler2003}
E.~Zitzler, L.~Thiele, M.~Laumanns, C.~M. Fonseca, and V.~G. da Fonseca.
Performance assessment of multiobjective optimizers: An analysis and review.
\emph{IEEE Transactions on Evolutionary Computation}, 7(2):117--132, 2003.

\bibitem{deb2014}
K.~Deb and H.~Jain.
An evolutionary many-objective optimization algorithm using reference-point-based nondominated sorting approach, part {I}.
\emph{IEEE Transactions on Evolutionary Computation}, 18(4):577--601, 2014.

\bibitem{gong2019}
X.~Gong, T.~De~Pessemier, L.~Martens, and W.~Joseph.
Energy- and labor-aware flexible job shop scheduling under dynamic electricity pricing: A many-objective optimization investigation.
\emph{Journal of Cleaner Production}, 209:1077--1094, 2019.

\bibitem{mouzon2007}
G.~Mouzon, M.~B. Yildirim, and J.~Twomey.
Operational methods for minimization of energy consumption of manufacturing equipment.
\emph{International Journal of Production Research}, 45(18--19):4247--4271, 2007.

\bibitem{dai2019}
M.~Dai, D.~Tang, A.~Giret, and M.~A. Salido.
Multi-objective optimization for energy-efficient flexible job shop scheduling problem with transportation constraints.
\emph{Robotics and Computer-Integrated Manufacturing}, 59:143--157, 2019.

\bibitem{allahverdi2008}
A.~Allahverdi, C.~T. Ng, T.~C.~E. Cheng, and M.~Y. Kovalyov.
A survey of scheduling problems with setup times or costs.
\emph{European Journal of Operational Research}, 187(3):985--1032, 2008.

\bibitem{allahverdi2015}
A.~Allahverdi.
The third comprehensive survey on scheduling problems with setup times/costs.
\emph{European Journal of Operational Research}, 246(2):345--378, 2015.

\bibitem{sutton2018}
R.~S. Sutton and A.~G. Barto.
\emph{Reinforcement Learning: An Introduction}.
MIT Press, 2nd edition, 2018.

\bibitem{mnih2015}
V.~Mnih, K.~Kavukcuoglu, D.~Silver, A.~A. Rusu, J.~Veness, M.~G. Bellemare, A.~Graves, M.~Riedmiller, A.~K. Fidjeland, G.~Ostrovski, S.~Petersen, C.~Beattie, A.~Sadik, I.~Antonoglou, H.~King, D.~Kumaran, D.~Wierstra, S.~Legg, and D.~Hassabis.
Human-level control through deep reinforcement learning.
\emph{Nature}, 518:529--533, 2015.

\bibitem{schulman2017}
J.~Schulman, F.~Wolski, P.~Dhariwal, A.~Radford, and O.~Klimov.
Proximal policy optimization algorithms.
arXiv:1707.06347, 2017.

\bibitem{haarnoja2018}
T.~Haarnoja, A.~Zhou, P.~Abbeel, and S.~Levine.
Soft actor-critic: Off-policy maximum entropy deep reinforcement learning with a stochastic actor.
In \emph{Proceedings of the 35th International Conference on Machine Learning}, pages 1861--1870, 2018.

\bibitem{zhang2020}
C.~Zhang, W.~Song, Z.~Cao, J.~Zhang, P.~S. Tan, and C.~Xu.
Learning to dispatch for job shop scheduling via deep reinforcement learning.
In \emph{Advances in Neural Information Processing Systems}, volume~33, pages 1621--1632, 2020.

\bibitem{park2021}
J.~Park, J.~Chun, S.~H. Kim, Y.~Kim, and J.~Park.
Learning to schedule job-shop problems: representation and policy learning using graph neural network and reinforcement learning.
\emph{International Journal of Production Research}, 59(11):3360--3377, 2021.

\bibitem{luo2020}
S.~Luo.
Dynamic scheduling for flexible job shop with new job insertions by deep reinforcement learning.
\emph{Applied Soft Computing}, 91:106208, 2020.

\bibitem{lei2022}
K.~Lei, P.~Guo, W.~Zhao, Y.~Wang, L.~Qian, X.~Meng, and L.~Tang.
A multi-action deep reinforcement learning framework for flexible job-shop scheduling problem.
\emph{Expert Systems with Applications}, 205:117796, 2022.

\bibitem{song2023}
W.~Song, X.~Chen, Q.~Li, and Z.~Cao.
Flexible job-shop scheduling via graph neural network and deep reinforcement learning.
\emph{IEEE Transactions on Industrial Informatics}, 19(2):1600--1610, 2023.

\bibitem{wang2025}
R.~Wang, Y.~Jing, C.~Gu, S.~He, and J.~Chen.
End-to-end multitarget flexible job shop scheduling with deep reinforcement learning.
\emph{IEEE Internet of Things Journal}, 12(4):4420--4434, 2025.

\bibitem{zeng2025}
D.~Zeng, W.~Yang, Y.~Du, X.~Yang, Y.~Wang, and Y.~Lv.
Dynamic multi-objective flexible job-shop scheduling via dynamic dual attention network-based deep reinforcement learning.
\emph{The Journal of Supercomputing}, 2025.
doi:10.1007/s11227-025-07586-z.

\bibitem{li2025}
J.~Li, S.~Li, P.~He, and H.~Li.
A multi objective collaborative reinforcement learning algorithm for flexible job shop scheduling.
\emph{Scientific Reports}, 2025.
doi:10.1038/s41598-025-03681-6.

\bibitem{vaswani2017}
A.~Vaswani, N.~Shazeer, N.~Parmar, J.~Uszkoreit, L.~Jones, A.~N. Gomez, L.~Kaiser, and I.~Polosukhin.
Attention is all you need.
In \emph{Advances in Neural Information Processing Systems}, volume~30, 2017.

\bibitem{velickovic2018}
P.~Velickovic, G.~Cucurull, A.~Casanova, A.~Romero, P.~Li\`o, and Y.~Bengio.
Graph attention networks.
In \emph{International Conference on Learning Representations}, 2018.

\bibitem{kipf2017}
T.~N. Kipf and M.~Welling.
Semi-supervised classification with graph convolutional networks.
In \emph{International Conference on Learning Representations}, 2017.

\bibitem{hamilton2017}
W.~L. Hamilton, R.~Ying, and J.~Leskovec.
Inductive representation learning on large graphs.
In \emph{Advances in Neural Information Processing Systems}, volume~30, 2017.

\bibitem{kool2019}
W.~Kool, H.~van Hoof, and M.~Welling.
Attention, learn to solve routing problems!
In \emph{International Conference on Learning Representations}, 2019.

\bibitem{chen2024}
F.~Chen, W.~Xie, J.~Ma, J.~Chen, and X.~Wang.
Textile flexible job-shop scheduling based on a modified ant colony optimization algorithm.
\emph{Applied Sciences}, 14(10):4082, 2024.

\bibitem{xu2011}
J.~Xu, X.~Xu, and S.~Q. Xie.
Recent developments in dual resource constrained ({DRC}) system research.
\emph{European Journal of Operational Research}, 215(2):309--318, 2011.

\bibitem{flapper2002}
S.~D.~P. Flapper, J.~C. Fransoo, R.~A.~C.~M. Broekmeulen, and K.~Inderfurth.
Planning and control of rework in the process industries: A review.
\emph{Production Planning and Control}, 13(1):26--34, 2002.

\bibitem{iec62264}
International Electrotechnical Commission.
IEC 62264: Enterprise-control system integration.
Geneva. Standard family built on ISA-95.

\bibitem{iec62541}
International Electrotechnical Commission.
IEC 62541: OPC unified architecture.
Geneva.

\bibitem{lee2015}
J.~Lee, B.~Bagheri, and H.-A. Kao.
A cyber-physical systems architecture for {I}ndustry 4.0-based manufacturing systems.
\emph{Manufacturing Letters}, 3:18--23, 2015.

\bibitem{monostori2016}
L.~Monostori, B.~K\'ad\'ar, T.~Bauernhansl, S.~Kondoh, S.~Kumara, G.~Reinhart, O.~Sauer, G.~Schuh, W.~Sihn, and K.~Ueda.
Cyber-physical systems in manufacturing.
\emph{CIRP Annals}, 65(2):621--641, 2016.

\bibitem{xu2018}
L.~D. Xu, E.~L. Xu, and L.~Li.
Industry 4.0: state of the art and future trends.
\emph{International Journal of Production Research}, 56(8):2941--2962, 2018.

\bibitem{koren2010}
Y.~Koren and M.~Shpitalni.
Design of reconfigurable manufacturing systems.
\emph{Journal of Manufacturing Systems}, 29(4):130--141, 2010.

\bibitem{grieves2017}
M.~Grieves and J.~Vickers.
Digital twin: Mitigating unpredictable, undesirable emergent behavior in complex systems.
In F.-J. Kahlen, S.~Flumerfelt, and A.~Alves, editors, \emph{Transdisciplinary Perspectives on Complex Systems}, pages 85--113. Springer, 2017.

\bibitem{tao2019}
F.~Tao, H.~Zhang, A.~Liu, and A.~Y.~C. Nee.
Digital twin in industry: State-of-the-art.
\emph{IEEE Transactions on Industrial Informatics}, 15(4):2405--2415, 2019.

\bibitem{qi2018}
Q.~Qi and F.~Tao.
Digital twin and big data towards smart manufacturing and {I}ndustry 4.0: 360 degree comparison.
\emph{IEEE Access}, 6:3585--3593, 2018.

\bibitem{negri2017}
E.~Negri, L.~Fumagalli, and M.~Macchi.
A review of the roles of digital twin in {CPS}-based production systems.
\emph{Procedia Manufacturing}, 11:939--948, 2017.

\bibitem{cimino2019}
C.~Cimino, E.~Negri, and L.~Fumagalli.
Review of digital twin applications in manufacturing.
\emph{Computers in Industry}, 113:103130, 2019.

\bibitem{fuller2020}
A.~Fuller, Z.~Fan, C.~Day, and C.~Barlow.
Digital twin: Enabling technologies, challenges and open research.
\emph{IEEE Access}, 8:108952--108971, 2020.

\bibitem{lu2020}
Y.~Lu, C.~Liu, K.~I.-K. Wang, H.~Huang, and X.~Xu.
Digital twin-driven smart manufacturing: Connotation, reference model, applications and research issues.
\emph{Robotics and Computer-Integrated Manufacturing}, 61:101837, 2020.

\bibitem{si2011}
X.-S. Si, W.~Wang, C.-H. Hu, and D.-H. Zhou.
Remaining useful life estimation --- a review on the statistical data driven approaches.
\emph{European Journal of Operational Research}, 213(1):1--14, 2011.

\bibitem{si2012}
X.-S. Si, W.~Wang, C.-H. Hu, D.-H. Zhou, and M.~G. Pecht.
Remaining useful life estimation based on a nonlinear diffusion degradation process.
\emph{IEEE Transactions on Reliability}, 61(1):50--67, 2012.

\bibitem{lei2018}
Y.~Lei, N.~Li, L.~Guo, N.~Li, T.~Yan, and J.~Lin.
Machinery health prognostics: A systematic review from data acquisition to {RUL} prediction.
\emph{Mechanical Systems and Signal Processing}, 104:799--834, 2018.

\bibitem{gebraeel2005}
N.~Z. Gebraeel, M.~A. Lawley, R.~Li, and J.~K. Ryan.
Residual-life distributions from component degradation signals: A {B}ayesian approach.
\emph{IIE Transactions}, 37(6):543--557, 2005.

\bibitem{kalman1960}
R.~E. Kalman.
A new approach to linear filtering and prediction problems.
\emph{Journal of Basic Engineering}, 82(1):35--45, 1960.

\bibitem{arulampalam2002}
M.~S. Arulampalam, S.~Maskell, N.~Gordon, and T.~Clapp.
A tutorial on particle filters for online nonlinear/non-{G}aussian {B}ayesian tracking.
\emph{IEEE Transactions on Signal Processing}, 50(2):174--188, 2002.

\bibitem{charnes1959}
A.~Charnes and W.~W. Cooper.
Chance-constrained programming.
\emph{Management Science}, 6(1):73--79, 1959.

\bibitem{birge2011}
J.~R. Birge and F.~Louveaux.
\emph{Introduction to Stochastic Programming}.
Springer, 2nd edition, 2011.

\bibitem{nguyen2019}
K.~T.~P. Nguyen and K.~Medjaher.
A new dynamic predictive maintenance framework using deep learning for failure prognostics.
\emph{Reliability Engineering and System Safety}, 188:251--262, 2019.

\bibitem{wang2024cie}
L.~Wang, B.~Li, and X.~Zhao.
Multi-objective predictive maintenance scheduling models integrating remaining useful life prediction and maintenance decisions.
\emph{Computers and Industrial Engineering}, 2024.

\bibitem{cbm2024}
L.~Li, L.~Deng, B.~Tang, and F.~Wang.
Bi-objective flexible job shop scheduling on machines considering condition-based maintenance activities.
\emph{Proceedings of the Institution of Mechanical Engineers, Part O: Journal of Risk and Reliability}, 238(6):1244--1255, 2024.

\bibitem{pandhare2024}
V.~Pandhare, E.~Negri, L.~Ragazzini, L.~Cattaneo, M.~Macchi, and J.~Lee.
Digital twin-enabled robust production scheduling for equipment in degraded state.
2024.

\bibitem{zhu2024pdm}
X.~Zhu, H.~Wang, and T.~Yuan.
An online approach for joint optimization of data-driven predictive maintenance and production planning.
In \emph{Proceedings of the European Safety and Reliability Conference (ESREL)}, 2024.

\end{thebibliography}


\end{document}
